% ============================================================
%  FinanceBench — Paper completo para Overleaf
%  Compilar con: pdflatex -> bibtex -> pdflatex -> pdflatex
%  Subir a Overleaf: este archivo + references.bib + carpeta figures/
% ============================================================
\documentclass[12pt,a4paper]{article}

% --- Codificación y tipografía ---
\usepackage[utf8]{inputenc}
\usepackage[T1]{fontenc}
\usepackage[spanish,es-tabla]{babel}
\usepackage{lmodern}
\usepackage{microtype}

% --- Márgenes ---
\usepackage[top=2.5cm,bottom=2.5cm,left=3cm,right=2.5cm]{geometry}

% --- Interlineado ---
\usepackage{setspace}
\onehalfspacing

% --- Matemáticas ---
\usepackage{amsmath}
\usepackage{amssymb}

% --- Gráficos ---
\usepackage{graphicx}
\graphicspath{{figures/}}
\usepackage{float}
\usepackage{caption}
\captionsetup{labelfont=bf, font=small, format=plain, justification=justified}

% --- Tablas ---
\usepackage{booktabs}
\usepackage{tabularx}
\usepackage{longtable}
\usepackage{multirow}
\usepackage{array}
\usepackage{xcolor}
\usepackage{colortbl}

% --- Bibliografía ---
\usepackage[round,authoryear,sort]{natbib}
\bibliographystyle{apalike}

% --- Hipervínculos ---
\usepackage[hidelinks,colorlinks=false]{hyperref}

% --- Párrafos ---
\setlength{\parindent}{0.5cm}
\setlength{\parskip}{4pt}

% --- Encabezado personalizado ---
\usepackage{fancyhdr}
\pagestyle{fancy}
\fancyhf{}
\fancyhead[L]{\small\textit{FinanceBench}}
\fancyhead[R]{\small\thepage}
\renewcommand{\headrulewidth}{0.4pt}

% ============================================================
\begin{document}

% --- Portada ---
\begin{titlepage}
  \centering
  \vspace*{2cm}
  {\LARGE\bfseries FinanceBench: Un Benchmark Integral para Evaluar \\[0.4em]
    Modelos de Lenguaje Grande en Tareas de \\[0.4em]
    Análisis Financiero\par}
  \vspace{1.5cm}
  {\large [afiliación institucional omitida]\\
    Cartagena, Colombia\par}
  \vspace{0.8cm}
  {\normalsize Línea de Investigación: Inteligencia Artificial Aplicada\\
    para la Democratización del Conocimiento Financiero\par}
  \vspace{1.5cm}
  {\large 2025\par}
  \vspace{2cm}
  \noindent\rule{\textwidth}{0.4pt}
  \begin{abstract}
    \noindent Los Modelos de Lenguaje Grande (LLM) se han convertido en herramientas de acceso cotidiano para análisis financiero, pero hasta la fecha no existía un benchmark que evaluara sistemáticamente su capacidad para realizar cálculos cuantitativos de viabilidad. Presentamos FinanceBench, el primer benchmark integral de análisis financiero cuantitativo: 1.000 casos de prueba distribuidos en 10 categorías de tareas financieras, validados por un panel de 5 expertos con coeficiente kappa de Cohen $\kappa = 0{,}87$, con valores oráculo calculados determinísticamente bajo la convención de flujos de caja para el accionista (FCFE). Evaluamos 9 LLM de frontera ---Claude Opus 4.6, Claude Sonnet 4.6, GPT-5.4, Gemini 3.1 Pro, GPT-5.3, Claude Haiku 4, DeepSeek Chat, Grok Code y Gemini 3 Flash--- bajo condiciones estandarizadas (prompt zero-shot, temperatura 0, 9.000 evaluaciones de línea base). Los resultados revelan una jerarquía de tres niveles: nivel superior ($>81\%$ de exactitud, 4 modelos), nivel intermedio ($69$--$79\%$, 3 modelos) y nivel bajo ($<63\%$, 2 modelos), con una brecha total de 30,83 puntos porcentuales entre el primero (Claude Opus 4.6, 88,47\%) y el último (Gemini 3 Flash, 57,64\%). Identificamos una ``hipótesis de la profundidad computacional'' que explica el gradiente de dificultad: la exactitud degrada de manera aproximadamente logarítmica con el número de operaciones dependientes, desde más del 97\% en depreciación hasta menos del 79\% en VPN y TIR. La composición de errores cambia cualitativamente con el nivel del modelo: los de nivel superior cometen principalmente errores aritméticos (47,2\%), mientras que los de nivel inferior acumulan errores conceptuales (31,8\%) y alucinaciones (19,4\%). FinanceBench, su dataset y código se publican como recurso de acceso abierto para la comunidad investigadora.
  \end{abstract}
  \noindent\rule{\textwidth}{0.4pt}
  \vspace{1cm}
  \noindent\textbf{Palabras clave:} grandes modelos de lenguaje, benchmarks de IA, análisis financiero cuantitativo, evaluación de LLM, democratización financiera, razonamiento numérico, VPN, TIR, América Latina.
\end{titlepage}

\tableofcontents
\newpage

% ============================================================
\section{Introducción}
% ============================================================

Pocas transformaciones tecnológicas recientes han alterado tan rápidamente la forma en que las personas interactúan con la información como la irrupción de los Large Language Models (LLM, por sus siglas en inglés; en español, Modelos de Lenguaje Grande). Un LLM es, en esencia, un sistema de inteligencia artificial entrenado con cantidades enormes de texto ---libros, artículos científicos, páginas web, foros, documentos técnicos--- que ha aprendido a identificar patrones lingüísticos lo suficientemente profundos como para generar respuestas coherentes, redactar documentos complejos y, en muchos casos, resolver problemas que hasta hace poco se consideraban exclusivos del intelecto humano. A diferencia de un motor de búsqueda, que localiza información existente y la presenta al usuario, un LLM puede razonar sobre esa información, combinar ideas de fuentes distintas y formular respuestas originales que no estaban previamente escritas en ningún sitio. Los sistemas más conocidos ---GPT-4 y su sucesor GPT-5 de OpenAI, Claude de Anthropic, Gemini de Google y LLaMA de Meta--- han pasado de ser curiosidades técnicas de laboratorio a herramientas de uso cotidiano en menos de tres años \citep{Brown2020, Touvron2023}. Según estimaciones recientes, el 71\% de las organizaciones encuestadas a nivel global reportan usar IA generativa de forma regular en al menos una función de negocio, frente al 65\% apenas un año antes \citep{McKinsey2025}, y la adopción en sectores como la salud, la educación, el derecho y las finanzas crece a un ritmo que supera las previsiones más optimistas de 2022. El 2025 AI Index Report de Stanford HAI confirma que el 78\% de las organizaciones usan IA en al menos una función operativa \citep{StanfordHAI2025}. La escala de este fenómeno no tiene precedente en la historia reciente de la inteligencia artificial: mientras que tecnologías anteriores como los asistentes de voz o los sistemas de recomendación tardaron una década en alcanzar penetración masiva, los LLM han logrado una adopción global en una fracción de ese tiempo.

El sector financiero se encuentra entre los que más rápidamente han abrazado esta nueva generación de herramientas. Y la razón es intuitiva: el análisis financiero ---la disciplina que permite determinar si un proyecto de inversión es viable, cuánto vale hoy el dinero que se recibirá mañana, o cuántas unidades necesita vender una empresa para cubrir sus costos--- es, al mismo tiempo, una actividad intelectualmente exigente y profundamente necesaria para millones de emprendedores, pequeñas empresas y tomadores de decisiones que carecen de acceso a asesoría profesional. En América Latina, esta tensión es particularmente aguda. El ecosistema emprendedor de la región ha crecido de manera sostenida durante la última década: solo en Colombia, la creación anual de nuevas empresas pasó de 278.302 en 2020 a un pico de 310.731 en 2022, impulsada por plataformas digitales, programas de apoyo gubernamental y una creciente cultura de innovación, aunque con una moderación posterior que sitúa la cifra de 2024 en 297.475 nuevas empresas \citep{Confecamaras2025}. Sin embargo, la mayoría de estos emprendedores no cuenta con formación formal en finanzas ni con los recursos para contratar un analista financiero profesional, cuyo costo puede superar varias veces el presupuesto mensual de una microempresa. En ese contexto, la promesa de un LLM que pueda calcular la Tasa Interna de Retorno de un proyecto, proyectar flujos de caja a cinco años o evaluar la viabilidad de un plan de negocio resulta irresistible. Según el informe del World Economic Forum y el Cambridge Centre for Alternative Finance, el 80\% de las empresas FinTech ---aquellas que combinan tecnología con servicios financieros--- están implementando inteligencia artificial en múltiples dominios de negocio \citep{WEFCCAF2025}, y un estudio de NVIDIA revela que el 46\% de las instituciones financieras ya usan LLM en sus operaciones, con un 55\% adicional explorando activamente la tecnología \citep{NVIDIA2024}. Emprendedores en Cartagena, Medellín o Ciudad de México ya emplean ChatGPT para redactar proyecciones financieras o evaluar estructuras de capital, y plataformas educativas integran asistentes de IA para enseñar conceptos como el Valor Presente Neto o el punto de equilibrio. La pregunta no es si los LLM se usarán para análisis financiero ---ya se están usando masivamente---, sino si producen resultados lo suficientemente confiables como para fundamentar decisiones con consecuencias económicas reales.

Y es aquí donde aparece el problema central que motiva este estudio. Un benchmark ---un conjunto estandarizado de pruebas que permite medir y comparar el rendimiento de diferentes sistemas bajo las mismas condiciones, de manera reproducible, como un examen estandarizado que todos los estudiantes presentan en igualdad de circunstancias--- es el instrumento fundamental que la ciencia de la computación utiliza para responder exactamente ese tipo de preguntas. Los benchmarks han sido el motor de progreso en casi todas las áreas de la inteligencia artificial: ImageNet catalizó la revolución del reconocimiento de imágenes, GLUE y SuperGLUE establecieron los estándares para comprensión de lenguaje natural \citep{Wang2019}, y MMLU permitió medir el conocimiento general de los LLM en 57 disciplinas \citep{Hendrycks2021}. Existen benchmarks de alta calidad para evaluar la capacidad matemática de estos modelos \citep{Cobbe2021, Hendrycks2021b}, para medir su desempeño en tareas médicas (MedQA), jurídicas (LegalBench) e incluso de programación (HumanEval). Sin embargo, no existe un benchmark diseñado específicamente para evaluar LLM en las tareas de análisis financiero que emprendedores y analistas realizan cotidianamente: cálculo de Valor Presente Neto, estimación de la Tasa Interna de Retorno, proyección de flujos de caja, análisis de punto de equilibrio, evaluación de estructuras de capital o generación de escenarios económicos. Los benchmarks generales como MMLU incluyen algunas preguntas de finanzas, pero se limitan a preguntas conceptuales de opción múltiple que no exigen cálculo numérico ni razonamiento de múltiples pasos; es la diferencia entre saber que el Valor Presente Neto existe y ser capaz de calcularlo correctamente con datos reales. Evaluaciones recientes como FinBen \citep{Xie2024} y PIXIU \citep{Xie2023} han avanzado en la evaluación de LLM en tareas de procesamiento de lenguaje natural financiero ---análisis de sentimiento, clasificación de noticias, extracción de entidades---, pero no cubren el tipo de razonamiento cuantitativo estructurado que define el análisis de viabilidad financiera. Las consecuencias prácticas de esta brecha evaluativa son serias: un desarrollador de FinTech que necesita elegir qué modelo integrar en su plataforma de asesoría financiera carece de datos comparativos rigurosos para tomar esa decisión; un emprendedor que recibe una proyección de flujo de caja generada por IA no tiene forma de saber si el modelo aplicó la fórmula correcta o si inventó datos inexistentes; y un investigador que quiere publicar un estudio sobre capacidades financieras de LLM no cuenta con un marco de referencia reproducible contra el cual comparar sus resultados.

El análisis financiero plantea un desafío particularmente difícil para los Modelos de Lenguaje Grande, y comprender por qué requiere entender cómo funcionan estos sistemas por dentro. Un LLM genera texto token a token ---cada token es una unidad de texto, típicamente una palabra o parte de una palabra--- prediciendo cuál es la siguiente pieza más probable dada la secuencia anterior. Este mecanismo, extraordinariamente eficaz para producir prosa coherente, no fue diseñado para realizar aritmética precisa. Cuando un modelo calcula el Valor Presente Neto de un proyecto a cinco años, debe encadenar múltiples operaciones de descuento, cada una dependiente del resultado anterior; un error de redondeo o de signo en el primer año se propaga y potencialmente se amplifica a lo largo de toda la cadena. La literatura ha documentado extensamente estas limitaciones: \citet{Wei2022} demostraron que técnicas como el razonamiento de cadena de pensamiento (\textit{Chain of Thought}), en el que se solicita al modelo que explique su proceso paso a paso antes de dar la respuesta final, mejoran sustancialmente el rendimiento en problemas aritméticos, pero no eliminan los errores. \citet{Li2024} mostraron en NumericBench que incluso modelos de última generación como GPT-4o fallan sistemáticamente en tareas numéricas aparentemente simples cuando se evalúan bajo condiciones controladas. A estas limitaciones aritméticas se suma un fenómeno especialmente peligroso en el contexto financiero: las alucinaciones, situaciones en las que el modelo genera información incorrecta o completamente inventada ---una tasa de interés que no estaba en los datos, un flujo de caja que no corresponde a ninguna proyección razonable--- con la misma confianza y fluidez con la que presenta información correcta, sin ofrecer al usuario ninguna señal de que algo está mal \citep{Huang2024}. En finanzas, una alucinación no es una anécdota curiosa: es un número falso que puede fundamentar una decisión de inversión equivocada.

Frente a esta problemática, presentamos FinanceBench, un benchmark integral diseñado para evaluar de manera sistemática y reproducible la capacidad de los Modelos de Lenguaje Grande en tareas fundamentales de análisis financiero. FinanceBench se compone de 1.000 casos de prueba distribuidos equitativamente en 10 categorías de tareas ---Cálculo de Valor Presente Neto, Cálculo de Tasa Interna de Retorno, Proyección de Flujo de Caja, Generación de Escenarios, Análisis de Punto de Equilibrio, Evaluación de Estructura de Capital, Cálculo de Indicadores Financieros, Interpretación de Estados Financieros, Análisis de Sensibilidad y Recomendaciones de Viabilidad---, cada uno con una solución correcta verificada por un panel de contadores certificados, analistas financieros y docentes universitarios de finanzas. Los casos están estratificados en tres niveles de dificultad (simple, medio y complejo) y reflejan escenarios realistas de emprendimientos y pequeñas empresas en el contexto colombiano y latinoamericano. Evaluamos 9 LLM ---incluyendo modelos propietarios como Claude Opus 4.6, Claude Sonnet 4.6, GPT-5.4, Gemini 3.1 Pro, GPT-5.3, Claude Haiku 4 y Gemini 3 Flash, así como modelos de código abierto como DeepSeek Chat y Grok Code--- bajo condiciones idénticas de prompting (la instrucción que se le da al modelo), generando aproximadamente 2,66 millones de comparaciones campo a campo contra valores oráculo. Los resultados revelan hallazgos que anticipamos aquí sin desarrollar completamente: existe una jerarquía clara de tres niveles de rendimiento, con una brecha de más de 30 puntos porcentuales entre el mejor y el peor modelo; las tareas que requieren cálculo iterativo (como el VPN y la TIR) son consistentemente más difíciles que las que involucran aritmética simple; y los tipos de error cambian cualitativamente a medida que mejora la capacidad del modelo, pasando de errores conceptuales ---aplicar la fórmula equivocada--- a errores aritméticos menores. El dataset completo y el código de evaluación se publican como recurso abierto.

La pertinencia de este estudio se amplifica cuando se lo sitúa en el contexto latinoamericano. La gran mayoría de los benchmarks existentes para evaluación de LLM han sido desarrollados en inglés, por equipos de investigación en Estados Unidos o Europa, y reflejan supuestos económicos, regulatorios y culturales que no necesariamente se trasladan a las economías emergentes de la región. Un emprendedor en Cartagena que evalúa la viabilidad de una cafetería no enfrenta los mismos costos de operación, tasas de descuento, estructuras impositivas ni condiciones de mercado que uno en San Francisco. La brecha de acceso a asesoría financiera profesional es, además, considerablemente más profunda en América Latina: según datos de la Asociación Nacional de Instituciones Financieras, menos del 25\% de las MiPyMEs en Colombia accede al crédito formal \citep{ANIF2024}, el 69\% de las microempresas no lleva contabilidad, y la tasa de supervivencia empresarial a cinco años es de apenas el 33,5\%, frente a tasas superiores al 40--50\% en el promedio de países de la OCDE \citep{Confecamaras2023}. En este contexto, los LLM representan una vía potencialmente democratizadora ---accesible, escalable, disponible las 24 horas---, pero solo si sus resultados son confiables. FinanceBench busca proporcionar exactamente la evidencia que permite determinar cuándo lo son y cuándo no.

% ============================================================
\section{Planteamiento del Problema}\label{sec:planteamiento}
% ============================================================

\subsection{Descripción del Problema}

La velocidad con la que los Modelos de Lenguaje Grande se han infiltrado en las prácticas financieras cotidianas ha generado un desajuste profundo entre uso y comprensión. Por un lado, una oleada de adopción sin precedentes: emprendedores que pegan sus datos de costos e ingresos en una ventana de chat y reciben, en segundos, una proyección de flujo de caja a cinco años con Valor Presente Neto incluido; startups de FinTech que integran APIs de modelos generativos para ofrecer ``asesoría financiera inteligente'' a miles de usuarios simultáneamente; instituciones educativas que incorporan chatbots como tutores de contabilidad y finanzas para estudiantes que nunca han abierto una hoja de cálculo. Por otro lado, una ausencia casi total de estándares que permitan saber si esos números que el modelo produce son correctos, aproximadamente correctos, o radicalmente equivocados.

Imaginemos un caso concreto para dimensionar el problema. María, una emprendedora en Cartagena, quiere abrir una cafetería. Ha estimado que necesita una inversión inicial de 50 millones de pesos colombianos y espera ingresos crecientes durante cinco años. Le pide a un LLM que calcule el Valor Presente Neto (VPN) de su proyecto ---es decir, que determine si el dinero que espera ganar en el futuro, una vez descontado el costo de esperar, justifica la inversión inicial---. El modelo le devuelve un número positivo y le dice que el proyecto es viable. María, que no tiene formación en finanzas y confía en la herramienta, procede a solicitar un crédito bancario basándose en esa proyección. Pero ¿qué pasa si el modelo aplicó mal la tasa de descuento? ¿Qué pasa si ``alucinó'' un flujo de caja en el tercer año que no corresponde a ningún dato ingresado? ¿Qué pasa si usó una fórmula de interés simple cuando debía usar interés compuesto? María no tiene forma de saberlo, porque no existe un marco de referencia público que le diga a ella ---ni al desarrollador de la herramienta, ni al regulador financiero, ni al investigador académico--- qué tan frecuentemente ocurren estos errores, en qué tipos de cálculos son más comunes, y qué modelos son más o menos confiables.

Este escenario no es hipotético: es la realidad cotidiana de millones de usuarios. La explosión de adopción de LLM en finanzas ha ocurrido en un vacío evaluativo. No existe un benchmark estandarizado y validado por expertos que permita comparar objetivamente la capacidad de diferentes modelos para realizar las tareas de análisis financiero que emprendedores, analistas y asesores ejecutan a diario. Los benchmarks generales de la comunidad de inteligencia artificial ---MMLU \citep{Hendrycks2021}, HELM \citep{Liang2023}, BIG-Bench \citep{Srivastava2022}--- evalúan conocimiento amplio y habilidades generales de razonamiento, pero no están diseñados para medir la precisión numérica en cálculos financieros de múltiples pasos ni la capacidad de interpretar resultados en contextos de negocio específicos. Los benchmarks financieros existentes, como FinBen \citep{Xie2024}, han avanzado significativamente en tareas de procesamiento de lenguaje natural financiero, pero su cobertura del razonamiento cuantitativo estructurado ---calcular un VPN, estimar una TIR, construir un punto de equilibrio--- es limitada. El resultado es un panorama donde las herramientas se adoptan masivamente antes de que la ciencia haya podido evaluar si funcionan.

La gravedad de esta situación se intensifica por un rasgo particular de los LLM: su capacidad para producir respuestas incorrectas con total fluidez y aparente confianza. A diferencia de una calculadora, que simplemente muestra ``ERROR'' cuando recibe una operación inválida, un LLM que comete un error financiero lo presenta envuelto en un párrafo bien redactado, con explicaciones detalladas que suenan perfectamente razonables, y sin ninguna advertencia de que el número central está equivocado. Este fenómeno, conocido como alucinación en la literatura de IA \citep{Huang2023}, es especialmente peligroso en finanzas porque las alucinaciones financieras son, por definición, indistinguibles de los análisis correctos para un usuario sin formación técnica. Un flujo de caja inventado por el modelo tiene exactamente la misma apariencia que uno calculado correctamente; la única diferencia está en los números internos, y detectarla requiere rehacer el cálculo completo de forma independiente.

\subsection{Manifestaciones Específicas del Problema}

El vacío evaluativo descrito se manifiesta de maneras concretas y observables que afectan a distintos actores del ecosistema.

En primer lugar, la ausencia de datos de referencia validados impide cualquier forma de comparación objetiva. Si un equipo de desarrollo quiere determinar si GPT-5.4 calcula el Valor Presente Neto con mayor precisión que Gemini 3.1 Pro, no tiene acceso a un conjunto estandarizado de problemas financieros con soluciones verificadas por expertos contra el cual medir ambos modelos en igualdad de condiciones. Cada equipo construye sus propias pruebas ad hoc, con criterios de evaluación distintos, y los resultados son incomparables entre estudios. Como señalan \citet{Liang2023} en el desarrollo de HELM, la estandarización de condiciones de evaluación ---mismos problemas, mismas métricas, mismos protocolos--- es un prerrequisito indispensable para cualquier conclusión válida sobre capacidades de modelos.

En segundo lugar, la selección no informada de modelos por parte de desarrolladores de aplicaciones financieras constituye una práctica generalizada con consecuencias potencialmente graves. Cuando una startup de FinTech elige integrar un modelo particular en su plataforma de asesoría financiera, esa decisión suele basarse en reputación de marca, costo de API, velocidad de respuesta o rendimiento en benchmarks generales como MMLU, ninguno de los cuales predice de manera confiable el rendimiento en cálculos financieros específicos. Nuestros propios resultados lo confirman de manera contundente: la brecha de 30,83 puntos porcentuales entre el mejor y el peor modelo de nuestro benchmark demuestra que la elección del modelo tiene consecuencias cuantitativamente enormes que no son visibles sin una evaluación especializada.

En tercer lugar, el riesgo para usuarios finales sin formación técnica es el aspecto más preocupante desde una perspectiva de impacto social. El emprendedor que utiliza un LLM para evaluar la viabilidad de su negocio no tiene, por lo general, la capacidad de verificar si el modelo aplicó correctamente la fórmula de descuento, si la tasa interna de retorno fue estimada con el método iterativo apropiado, o si los escenarios pesimistas reflejan adecuadamente los riesgos. La asimetría de información entre el usuario y la herramienta es total: el modelo ``sabe'' cosas que el usuario no puede verificar, y cuando el modelo se equivoca, el usuario no puede detectarlo. En economías emergentes como la colombiana, donde los errores en la planificación financiera de una microempresa pueden significar la pérdida de los ahorros de toda una familia, esta asimetría adquiere una dimensión ética ineludible.

En cuarto lugar, la imposibilidad de comparar estudios entre investigadores frena el progreso científico en el campo. Sin un benchmark compartido, cada grupo de investigación mide cosas distintas de maneras distintas, y la acumulación de conocimiento ---que es la base del método científico--- se vuelve imposible. Un investigador que reporta que ``GPT-4 obtiene un 82\% de precisión en tareas financieras'' no está diciendo algo comparable a otro que reporta que ``Claude obtiene un 79\%'', porque las tareas, las métricas y los criterios de corrección son probablemente diferentes. Esta fragmentación replica un problema que la comunidad de IA general ya resolvió con benchmarks como MMLU y HELM, pero que persiste en el dominio financiero.

Finalmente, existe un optimismo injustificado en las capacidades financieras de los LLM que se alimenta, paradójicamente, de la falta de evaluación rigurosa. En ausencia de datos empíricos que muestren dónde fallan estos sistemas, es fácil asumir que ``funcionan bien'' basándose en demostraciones anecdóticas o en su innegable capacidad para generar texto convincente sobre temas financieros. Nuestros datos cuentan una historia más matizada: incluso el mejor modelo evaluado (Claude Opus 4.6, con 88,47\% de exactitud) comete errores en más de una de cada diez evaluaciones, y para los modelos de nivel inferior, la tasa de error puede superar el 40\%. Estas cifras no son evidentes a simple vista cuando se interactúa casualmente con un chatbot, y solo emergen bajo una evaluación sistemática como la que FinanceBench proporciona.

\subsection{Justificación e Importancia del Estudio}

La relevancia de este estudio se fundamenta en tres dimensiones complementarias que, en conjunto, constituyen un argumento sólido a favor de su realización en el momento presente.

Desde una perspectiva teórica, FinanceBench contribuye al conocimiento científico sobre las capacidades y limitaciones de los Modelos de Lenguaje Grande en un dominio que combina razonamiento numérico, conocimiento de dominio específico e interpretación contextual. La literatura existente ha estudiado extensamente el razonamiento matemático de los LLM \citep{Wei2022, Cobbe2021, Hendrycks2021b}, pero el razonamiento financiero introduce complejidades adicionales que los benchmarks matemáticos no capturan: la necesidad de seleccionar la fórmula correcta entre varias posibles (¿VPN o TIR? ¿interés simple o compuesto?), de interpretar el resultado en un contexto de negocio (un VPN positivo no significa lo mismo para una cafetería que para una plataforma tecnológica), y de manejar múltiples escenarios económicos simultáneamente. Al evaluar 9 modelos en 1.000 casos con 10 tipos de tareas diferentes, FinanceBench genera una base empírica que permite responder preguntas teóricas que hasta ahora solo podían abordarse de manera especulativa: ¿escala el rendimiento financiero con el tamaño del modelo? ¿Los errores son aditivos o multiplicativos en cadenas de cálculo? ¿Existe un umbral de capacidad por debajo del cual el razonamiento financiero colapsa?

Desde una perspectiva práctica, los resultados del benchmark tienen aplicaciones inmediatas para al menos tres audiencias. Para los desarrolladores de aplicaciones de FinTech, proporcionan una guía basada en evidencia para seleccionar el modelo más adecuado según el tipo de tarea financiera y el nivel de precisión requerido; nuestros datos muestran, por ejemplo, que la diferencia entre usar un modelo de nivel superior y uno de nivel inferior puede representar más de 30 puntos porcentuales de exactitud, una brecha que se traduce directamente en la calidad del servicio ofrecido al usuario final. Para los emprendedores y usuarios finales, los resultados ofrecen por primera vez una referencia concreta sobre qué esperar de estas herramientas: en qué tareas pueden confiar, en cuáles deben verificar, y en cuáles no deberían confiar en absoluto sin supervisión profesional. Para los reguladores financieros, FinanceBench aporta evidencia cuantitativa que puede informar marcos regulatorios sobre el uso de IA en asesoría financiera, un tema que está emergiendo en agendas regulatorias globales, desde el AI Act de la Unión Europea hasta las guías de la Autoridad Monetaria de Singapur \citep{MAS2025} y los informes del Financial Stability Board \citep{FSB2024}.

Desde una perspectiva contextual, la urgencia de este estudio es especialmente aguda en América Latina. La región concentra una combinación de factores que amplifica tanto la oportunidad como el riesgo de los LLM en finanzas: un ecosistema emprendedor dinámico pero con limitado acceso a asesoría profesional \citep{Estefan2023}, una penetración creciente de smartphones y conectividad digital que facilita el acceso a herramientas de IA, y un contexto económico donde los márgenes de error en la planificación financiera de micro y pequeñas empresas son estrechos. En este entorno, un LLM que produce un análisis financiero erróneo no es simplemente un inconveniente técnico: puede ser la diferencia entre un emprendimiento que prospera y uno que fracasa llevándose consigo los ahorros de una familia. FinanceBench busca proporcionar la evidencia necesaria para que esta adopción ocurra de manera informada y responsable, no a ciegas.

La combinación de estas tres dimensiones ---avance teórico, utilidad práctica inmediata y urgencia contextual--- constituye, a nuestro juicio, una justificación suficiente para la realización de este estudio en el momento actual. La alternativa ---permitir que la adopción masiva de LLM en finanzas continúe sin evaluación rigurosa--- tiene un costo de oportunidad científico y un riesgo social que este trabajo busca reducir.

% ============================================================
\section{Objetivos de Investigación}\label{sec:objetivos}
% ============================================================

\subsection*{Objetivo General}

Diseñar un benchmark integral para la evaluación sistemática de Modelos de Lenguaje Grande en tareas de análisis financiero, que permita comparar objetivamente el rendimiento de los modelos existentes y establecer estándares de calidad para aplicaciones de inteligencia artificial en finanzas.

\subsection*{Objetivos Específicos}

\begin{enumerate}
  \item Diseñar una taxonomía de tareas financieras relevantes para el análisis de viabilidad de negocios.
  \item Establecer métricas de evaluación apropiadas para medir precisión calculatoria y calidad de razonamiento financiero.
  \item Construir un dataset de 1.000 casos de prueba con soluciones verificadas por expertos.
  \item Evaluar el rendimiento de 9 LLM en el benchmark construido.
  \item Clasificar los tipos de errores cometidos por los modelos en cinco categorías (aritmético, conceptual, alucinación, formato y campo nulo/ausente) para identificar patrones de fallo diferenciados por nivel de capacidad.
\end{enumerate}

% ============================================================
\section{Marco Teórico}\label{sec:marco}
% ============================================================

Esta sección construye los cimientos conceptuales del estudio. No es un glosario sino una narrativa que introduce los conceptos necesarios, los conecta entre sí y justifica las decisiones de diseño de FinanceBench.

\subsection{Inteligencia Artificial y Modelos de Lenguaje Grande}

La inteligencia artificial (IA) es una rama de la ciencia de la computación que desarrolla sistemas capaces de realizar tareas que normalmente requieren inteligencia humana: entender texto, reconocer imágenes, traducir idiomas o diagnosticar enfermedades. Su evolución puede resumirse en tres paradigmas. Los \textit{sistemas expertos} (1960--1980) funcionaban con reglas escritas a mano. El \textit{aprendizaje automático} (\textit{machine learning}) permitió que los sistemas descubrieran patrones a partir de datos, en lugar de ser programados para cada caso. Y el \textit{aprendizaje profundo} (\textit{deep learning}) utilizó redes neuronales de múltiples capas para capturar patrones cada vez más abstractos \citep{LeCun2015}.

Los LLM representan el hito más reciente de esta trayectoria: sistemas de aprendizaje profundo entrenados con cantidades masivas de texto que aprenden a predecir y generar lenguaje de manera coherente. Si un motor de búsqueda encuentra documentos existentes, un LLM puede razonar sobre múltiples fuentes y construir una respuesta original. El mecanismo técnico que los hizo posibles es la arquitectura \textit{Transformer} \citep{Vaswani2017}, cuya innovación central ---el mecanismo de \textit{atención}--- permite al modelo ponderar qué partes del contexto son más relevantes para la tarea en curso, capturando relaciones entre palabras lejanas sin procesamiento secuencial.

A partir de esta arquitectura, la comunidad ha producido modelos de escala creciente. GPT-3 \citep{Brown2020}, con 175 mil millones de \textit{parámetros} ---los valores numéricos internos que el modelo ajusta durante el entrenamiento, análogos al número de conexiones en un circuito neuronal---, mostró capacidades emergentes: podía realizar tareas para las que no había sido entrenado explícitamente. Los modelos más recientes ---GPT-5 de OpenAI \citep{OpenAI2024}, Claude de Anthropic \citep{Anthropic2025}, Gemini de Google \citep{GoogleDeepMind2024}, LLaMA de Meta \citep{Touvron2023}--- representan la frontera actual, con capacidades que incluyen razonamiento de múltiples pasos y, de manera creciente, cálculos numéricos complejos. Es esta última capacidad la que los posiciona como candidatos para tareas de análisis financiero y, simultáneamente, la que exige evaluación rigurosa.

\subsection{Capacidades y Limitaciones de los LLM en Razonamiento Numérico}

El \textit{razonamiento verbal} ---comprender argumentos, sintetizar información--- es el terreno natural de los LLM. Sin embargo, el \textit{razonamiento numérico} presenta un panorama más irregular. Los LLM no tienen una ``calculadora interna''; cuando multiplican dos números, predicen el resultado más probable basándose en patrones de texto, no ejecutan una operación aritmética convencional. Un modelo puede resolver correctamente ``247 × 13'' y fallar en ``2.473,58 × 13,07'' porque la segunda combinación es menos frecuente en sus datos de entrenamiento.

GSM8K \citep{Cobbe2021} muestra que incluso modelos de gran escala cometen errores en cadenas de más de cuatro operaciones. MATH \citep{Hendrycks2021b} revela que los modelos tienden a cometer errores locales de álgebra dentro de razonamientos globalmente correctos. DROP \citep{Dua2019} evalúa razonamiento numérico sobre textos extensos ---un formato que se asemeja al análisis de estados financieros---. Y NumericBench \citep{Li2024} muestra que modelos como GPT-4o fallan incluso en tareas que parecerían triviales, como identificar el valor máximo en una columna de datos.

El \textit{Chain of Thought} (CoT), propuesto por \citet{Wei2022}, mejora el rendimiento al pedir al modelo que explique su proceso paso a paso. La idea es que así como un estudiante que escribe cada paso comete menos errores que uno que hace todo mentalmente, un LLM que verbaliza sus pasos intermedios produce resultados más precisos. \citet{Lewkowycz2022} mostraron con Minerva que CoT combinado con votación por mayoría alcanza rendimientos cercanos al estado del arte en problemas STEM. Sin embargo, \citet{Lightman2023} demostraron que un modelo puede generar una cadena de pasos que suena lógica pero contiene un error numérico intermedio que invalida la conclusión, sin que la aparente solidez del razonamiento ofrezca señal alguna del fallo.

Las \textit{alucinaciones} constituyen la limitación más peligrosa en contexto financiero. Ocurren cuando el LLM genera información sin sustento en los datos proporcionados, presentándola con la misma fluidez que la información correcta. En finanzas, una alucinación puede ser una tasa de interés fabricada, un flujo de caja inventado, o una recomendación basada en números que el modelo generó de la nada. \citet{Huang2024} ofrecen una taxonomía detallada, y estudios posteriores confirman que la frecuencia de alucinaciones aumenta con la complejidad de la tarea. Nuestros resultados lo confirman: la tasa de alucinación oscila entre el 11,4\% de los errores (Claude Opus 4.6) y el 19,4\% (Gemini 3 Flash).

\subsection{Análisis Financiero: Conceptos Clave}

El análisis de viabilidad financiera involucra un conjunto de conceptos que el lector necesita comprender para interpretar los resultados del benchmark. Presentamos aquí los más relevantes, definidos en términos accesibles.

\textbf{Valor Presente Neto (VPN / NPV).} Un peso que recibirás dentro de un año no vale lo mismo que un peso que tienes hoy, porque el peso de hoy puede invertirse y generar rendimiento. El VPN toma todos los flujos de dinero futuros de un proyecto, los ``descuenta'' al presente usando una tasa que refleja el costo de esperar, y les resta la inversión. Si el resultado es positivo, el proyecto genera más valor del que consume. La diferencia entre calcular el VPN sobre los flujos totales del proyecto y sobre los flujos que le quedan al dueño después de pagar sus deudas es fundamental, como explicamos a continuación.

\textbf{Flujos de caja para el accionista (FCFE) vs.\ flujos para la firma (FCFF).} La diferencia entre estos dos conceptos es como la diferencia entre mirar cuánto gana un negocio en total (antes de pagar sus deudas) y cuánto le queda al dueño después de pagar el banco. Si María invirtió \$30 millones de su bolsillo y pidió \$20 millones prestados, lo que le interesa saber es cuánto le rinde \textit{su} inversión de \$30 millones, no la inversión total de \$50 millones. Los flujos para el accionista (FCFE) responden esa pregunta: son la utilidad neta más la depreciación menos el pago de capital de la deuda. FinanceBench adopta esta convención porque refleja la perspectiva del emprendedor que evalúa su inversión personal \citep{SapagChain2014, BerkDeMarzo2020}.

\textbf{Tasa Interna de Retorno (TIR / IRR).} Si el VPN responde ``¿este proyecto genera valor?'', la TIR responde ``¿cuánto rendimiento porcentual genera?''. Técnicamente, es la tasa de descuento que hace el VPN exactamente igual a cero. Calcularla requiere un proceso iterativo ---probar tasas hasta encontrar la correcta---, algo que para un LLM constituye un desafío significativo.

\textbf{Proyección de Flujo de Caja.} Estimar los ingresos y egresos de dinero a lo largo del tiempo, la base sobre la que se construyen todos los demás análisis.

\textbf{Generación de Escenarios.} Construir versiones alternativas del futuro ---optimista, normal, pesimista y controlado--- variando supuestos clave para evaluar cómo responde la viabilidad ante diferentes condiciones.

\textbf{Punto de Equilibrio.} ¿Cuántas unidades necesito vender para cubrir todos mis costos? La fórmula clásica: $Q = \text{costos fijos} / (\text{precio} - \text{costo variable por unidad})$.

\textbf{Estructura de Capital y WACC.} Cómo se financia una empresa (deuda vs.\ capital propio) y cuál es el costo promedio de esa mezcla. El WACC (\textit{Weighted Average Cost of Capital}) pondera cada fuente por su costo después de impuestos.

\textbf{Indicadores Financieros.} ROI (retorno sobre la inversión), ROE (retorno sobre el capital propio), margen bruto, EBITDA: métricas que resumen el desempeño de un negocio en cifras comparables.

\textbf{Interpretación de Estados Financieros.} Lectura y análisis del balance general, estado de resultados y estado de flujo de efectivo para comprender la salud financiera de una empresa.

\textbf{Análisis de Sensibilidad.} Identificar cuáles variables impactan más los resultados: ¿qué pasa si la tasa de descuento sube un punto? ¿Si los ingresos caen un 10\%?

\textbf{Recomendaciones de Viabilidad.} La síntesis de todos los análisis anteriores para emitir un juicio: ¿debe procederse con la inversión o no?

\subsection{Benchmarks en Inteligencia Artificial}

El concepto de benchmark es para la IA lo que el ensayo clínico es para la medicina: el mecanismo que permite determinar si un sistema funciona antes de desplegarlo. Un benchmark consta de tres componentes: un conjunto de tareas con soluciones conocidas, un protocolo de evaluación estandarizado, y métricas que cuantifican el rendimiento.

GLUE \citep{Wang2019} y SuperGLUE \citep{Wang2019b} establecieron el paradigma de evaluación multi-tarea con leaderboard público. MMLU \citep{Hendrycks2021} evaluó conocimiento en 57 áreas; su versión actualizada MMLU-Pro \citep{Yuan2024} corrigió problemas de contaminación de datos. HELM \citep{Liang2023} amplió el paradigma evaluando calibración, robustez, equidad y eficiencia. BIG-Bench \citep{Srivastava2022} incluyó tareas diseñadas para estar ``más allá'' de la capacidad actual.

Ninguno aborda adecuadamente el razonamiento financiero aplicado. MMLU pregunta qué \textit{es} el VPN, no pide \textit{calcularlo}. GSM8K y MATH evalúan aritmética sin fórmulas financieras. FinQA \citep{Chen2021} evalúa razonamiento numérico sobre reportes financieros pero no la generación de análisis completos. FinBen \citep{Xie2024}, el más comprehensivo, cubre 40+ datasets de NLP financiero pero no el cálculo estructurado de viabilidad. FinanceBench se posiciona en esa brecha.

\subsection{IA y Finanzas: Estado del Arte}

BloombergGPT \citep{Wu2023} demostró que el entrenamiento sobre un corpus financiero masivo mejora el rendimiento en NLP financiero. FinGPT \citep{Yang2023} propuso un framework de código abierto para fine-tuning de modelos con datos financieros. PIXIU \citep{Xie2023} combinó un modelo financiero, un dataset de instrucciones y un benchmark. Los surveys de \citet{Li2023} y \citet{Nie2024} sintetizan el campo y señalan que la evaluación de capacidades cuantitativas sigue subdesarrollada. Finance Agent Benchmark \citep{Bigeard2025} evalúa agentes LLM en tareas de equity research, un enfoque complementario al nuestro.

\subsection{Democratización del Conocimiento Financiero en América Latina}

En Colombia, las MiPyMEs representan el 99,4\% de las empresas formalizadas, de las cuales el 91,8\% son microempresas \citep{Asobancaria2022}. La tasa de supervivencia empresarial a cinco años es del 33,5\% ---solo un tercio de las empresas creadas en 2017 seguía activa en 2022---, frente a tasas superiores al 40--50\% en el promedio OCDE \citep{Confecamaras2023}. Menos del 25\% de las MiPyMEs accede al crédito formal, el 69\% no lleva contabilidad y el 39\% no separa las cuentas del hogar y del negocio \citep{Asobancaria2022}.

\citet{Estefan2023}, en un RCT con microemprendedores en Guatemala, demostraron que herramientas digitales mejoran conocimientos financieros y resultados económicos. \citet{Polanco2024}, con 2.440 estudiantes, encontró que un chatbot generativo mejoró puntajes de alfabetización financiera en 0,13 desviaciones estándar. \citet{Yang2025} reportó mejoras del 42\% en conocimiento financiero mediante IA generativa en regiones con recursos escasos.

Estas evidencias son alentadoras, pero llevan una advertencia: si los LLM que se utilizan para asesoría financiera producen resultados incorrectos, el daño potencial es mayor en las poblaciones más vulnerables. Como señalan \citet{CGAP2025} y \citet{FSB2024}, la promesa de inclusión financiera mediante IA solo se materializa si las herramientas son confiables. FinanceBench contribuye proporcionando la evidencia empírica para distinguir entre modelos confiables y modelos riesgosos.

% ============================================================
\section{Antecedentes de Investigación}\label{sec:antecedentes}
% ============================================================

La Sección~\ref{sec:marco} construyó el vocabulario conceptual; esta sección examina estudios previos, analiza lo que encontraron y, sobre todo, identifica lo que dejaron sin resolver. La conclusión es clara: pese al enorme volumen de investigación en la intersección de LLM y finanzas, la evaluación sistemática del razonamiento financiero cuantitativo sigue pendiente.

\subsection{Benchmarks Generales de LLM}

MMLU \citep{Hendrycks2021} evalúa conocimiento en 57 áreas mediante preguntas de opción múltiple de nivel profesional. Reveló que incluso los modelos más grandes obtenían puntuaciones inferiores al nivel humano experto, y que la confianza del modelo estaba frecuentemente mal calibrada ---se mostraba igual de seguro respondiendo correcta e incorrectamente---. Esta observación es directamente relevante para finanzas, donde un modelo que presenta un VPN erróneo con la misma confianza que uno correcto es peligroso. Sin embargo, las preguntas financieras de MMLU son exclusivamente conceptuales: nunca piden calcular un VPN a partir de datos numéricos. MMLU-Pro \citep{Yuan2024} elevó la dificultad pero mantuvo el formato conceptual.

HELM \citep{Liang2023} amplió el paradigma evaluando exactitud, calibración, robustez, equidad y eficiencia en más de 40 escenarios, estableciendo un estándar de transparencia reproducible que FinanceBench adopta. Sin embargo, las finanzas son, en terminología de sus autores, un ``escenario ausente''.

BIG-Bench \citep{Srivastava2022} reveló fenómenos como los avances súbitos de rendimiento al cruzar umbrales de escala ---tareas donde modelos pequeños fallan y modelos grandes aciertan, sin rendimiento intermedio---, un patrón análogo al que observamos en el punto de equilibrio.

La lección transversal: la estandarización rigurosa es indispensable, pero los benchmarks generales no pueden cubrir dominios que requieren conocimiento específico, formulación numérica y razonamiento multi-paso contextualizado.

\subsection{Benchmarks de Razonamiento Numérico y Matemático}

GSM8K \citep{Cobbe2021} introdujo el concepto de ``verificadores'' que evalúan si el razonamiento de un modelo es correcto, un principio afín a nuestro recálculo determinístico de valores oráculo. MATH \citep{Hendrycks2021b} documentó que los modelos cometen errores locales de aritmética dentro de razonamientos globalmente correctos ---exactamente el perfil de error que domina en los modelos de nivel superior de nuestro benchmark, donde el 47,2\% de los errores de Claude Opus 4.6 son aritméticos---. DROP \citep{Dua2019} evalúa razonamiento discreto sobre párrafos extensos, un formato que se asemeja al análisis de estados financieros. NumericBench \citep{Li2024} mostró que incluso modelos de última generación fallan en tareas numéricas aparentemente triviales.

La síntesis: los LLM han mejorado en matemáticas (GSM8K ha escalado del 57\% al $>$90\%), pero el razonamiento financiero involucra selección de fórmulas, interpretación contextual y manejo simultáneo de múltiples escenarios ---dimensiones que ningún benchmark matemático captura---.

\subsection{Benchmarks de Dominio Específico}

En medicina, MedQA evalúa preguntas del USMLE; en derecho, LegalBench evalúa razonamiento jurídico; en programación, HumanEval evalúa generación de código funcional. De estos precedentes extraemos tres lecciones: la taxonomía de tareas debe reflejar actividades reales del dominio, la validación por expertos del dominio es indispensable, y las métricas deben capturar múltiples dimensiones del rendimiento.

\subsection{Evaluaciones Previas de LLM en Tareas Financieras}

FinQA \citep{Chen2021} evalúa razonamiento numérico sobre tablas de reportes financieros corporativos, aportando un análisis de errores típicos que inspiró nuestro esquema de clasificación. Sin embargo, se circunscribe a preguntas sobre informes existentes; no evalúa la generación de análisis financieros completos.

FinBen \citep{Xie2024}, presentado en NeurIPS 2024, reúne más de 40 datasets y 20 tareas, evaluando más de 20 LLM. GPT-4 lidera en la mayoría de las tareas y los modelos de código abierto muestran mejora sostenida. FinanceBench se posiciona como complementario: mientras FinBen abarca amplitud (muchas tareas de NLP financiero), FinanceBench ofrece profundidad en cálculo cuantitativo estructurado con métricas numéricas finas.

PIXIU \citep{Xie2023} ofreció un ecosistema para LLM financieros pero se centra en clasificación y predicción. BloombergGPT \citep{Wu2023} mejoró NLP financiero pero su evaluación se limitó a tareas textuales y su naturaleza propietaria impide replicación.

Chatbot Arena \citep{Wang2024} introdujo comparaciones pareadas ciegas con jueces humanos y automatizados. Finance Agent Benchmark \citep{Bigeard2025} evalúa agentes LLM en tareas de equity research con acceso a herramientas externas ---complementario a FinanceBench, que aborda el componente numérico---.

\subsection{Síntesis y Brecha de Conocimiento}

\begin{table}[H]
\centering
\caption{Síntesis de antecedentes y brechas}
\label{tab:antecedentes}
\begin{tabularx}{\textwidth}{lXXX}
\toprule
\textbf{Estudio} & \textbf{Tareas evaluadas} & \textbf{Modelos} & \textbf{Limitación para nuestro propósito} \\
\midrule
MMLU \citep{Hendrycks2021} & 57 áreas, opción múltiple & $>$30 & Sin cálculo numérico financiero \\
HELM \citep{Liang2023} & 42 escenarios, multi-métrica & $>$30 & Finanzas ausentes \\
BIG-Bench \citep{Srivastava2022} & $>$200 tareas diversas & Múltiples & Sin foco financiero \\
GSM8K \citep{Cobbe2021} & Aritmética escolar & Múltiples & Sin fórmulas financieras \\
MATH \citep{Hendrycks2021b} & Matemáticas olímpicas & Múltiples & Sin dominio financiero \\
NumericBench \citep{Li2024} & 6 habilidades numéricas & Varios & Numeracia básica, no finanzas \\
FinQA \citep{Chen2021} & QA numérico sobre informes & Varios & Solo extracción, no generación \\
FinBen \citep{Xie2024} & 40+ datasets NLP financiero & 20+ & Énfasis NLP, no cálculo estructurado \\
PIXIU \citep{Xie2023} & 5 tareas NLP + predicción & Varios & Sin VPN, TIR, punto de equilibrio \\
BloombergGPT \citep{Wu2023} & NLP financiero & BloombergGPT & Propietario; sin cálculos \\
Finance Agent Bench.\ \citep{Bigeard2025} & Equity research agentic & Múltiples & Foco textual, no cálculo \\
Chatbot Arena \citep{Wang2024} & Conversación abierta & $>$50 & Sin finanzas ni métricas numéricas \\
\bottomrule
\end{tabularx}
\smallskip
\footnotesize\textit{Fuente: Elaboración propia.}
\end{table}

Tres brechas emergen. La \textbf{primera}: no existe un benchmark que evalúe sistemáticamente la capacidad de los LLM para realizar cálculos financieros validados por expertos. La \textbf{segunda} es metodológica: los estudios existentes usan muestras pequeñas, métricas limitadas y condiciones no estandarizadas. La \textbf{tercera} es contextual: toda la investigación refleja supuestos de economías anglosajonas, sin representación del contexto latinoamericano. FinanceBench llena estas tres brechas.

% ============================================================
\section{Metodología}\label{sec:metodologia}
% ============================================================

El criterio rector de esta sección es la reproducibilidad: cualquier investigador debe poder replicar exactamente este estudio siguiendo exclusivamente la descripción aquí presentada.

\subsection{Diseño General del Estudio}

Este trabajo constituye una evaluación comparativa sistemática mediante benchmark computacional. Diseñamos un conjunto estandarizado de pruebas financieras, lo aplicamos a 9 LLM bajo condiciones idénticas, y comparamos los resultados de manera objetiva. El diseño se inscribe dentro del paradigma de investigación cuantitativa.

El proyecto se ejecutó en tres fases. La \textbf{Fase 1: Construcción del Benchmark} (meses 1--3) abarcó diseño de taxonomía, generación de 1.000 casos, producción de soluciones determinísticas y validación por panel de expertos. La \textbf{Fase 2: Evaluación Baseline} (meses 4--5) consistió en la aplicación del benchmark a los 9 modelos, recolección de respuestas, extracción automatizada y cálculo de métricas. La \textbf{Fase 3: Análisis y Publicación} (mes 6) comprendió el análisis estadístico, clasificación de errores y preparación del dataset para publicación abierta. La separación en fases garantiza que la construcción del benchmark esté completa antes de que cualquier modelo sea evaluado.

\subsection{Principios de Diseño del Benchmark}

Cinco principios guiaron las decisiones de diseño. \textbf{Relevancia:} cada tarea representa una actividad de análisis financiero real. \textbf{Diversidad:} cobertura de razonamiento aritmético, conceptual, mixto e interpretativo. \textbf{Verificabilidad:} cada caso tiene respuesta correcta objetivamente determinable. \textbf{Realismo:} escenarios que reflejan condiciones plausibles de negocio en 10 sectores diferentes. \textbf{Gradación de dificultad:} distribución en simple (40\%), medio (40\%) y complejo (20\%).

\subsection{Convención de Flujos de Caja Adoptada}

Existen dos convenciones para calcular el VPN y la TIR: la del flujo de caja libre para la firma (FCFF), que se descuenta al WACC y evalúa el proyecto independientemente de cómo se financia; y la del flujo de caja libre para el accionista (FCFE), que se descuenta al costo de capital del accionista ($k_e$) y evalúa la rentabilidad desde la perspectiva del inversionista de capital propio. La diferencia es como mirar cuánto gana un negocio en total (antes de pagar deudas) versus cuánto le queda al dueño después de pagar al banco.

FinanceBench adopta la \textbf{convención FCFE}, formalizada en las siguientes ecuaciones:

\begin{align}
  \text{FCFE}_t &= \text{UN}_t + \text{Dep}_t - \text{PK}_t \label{eq:fcfe}\\
  \text{VPN} &= \sum_{t=1}^{T} \frac{\text{FCFE}_t}{(1+k_e)^t} - I_0 \label{eq:vpn}\\
  \text{TIR} &: \sum_{t=1}^{T} \frac{\text{FCFE}_t}{(1+\text{TIR})^t} - I_0 = 0 \label{eq:tir}
\end{align}

donde $\text{UN}_t$ es la utilidad neta del período $t$, $\text{Dep}_t$ la depreciación, $\text{PK}_t$ el pago de capital (amortización del principal) de la deuda, $k_e$ el costo de capital del accionista, $I_0$ la inversión inicial de capital propio, y $T$ el horizonte de evaluación.

Esta decisión se fundamenta en tres razones. Primera, el usuario objetivo ---emprendedores y analistas de PYMEs--- evalúa la rentabilidad desde la perspectiva del capital que aporta al proyecto \citep{SapagChain2014, BerkDeMarzo2020}. Segunda, mantiene consistencia interna: VPN y TIR usan la misma base de flujos. Tercera, es reproducible sin ambigüedad. La elección afecta la comparabilidad con benchmarks que usen WACC + FCFF, limitación reconocida en la Sección~\ref{sec:limitaciones}.

\subsection{Taxonomía de Tareas Financieras}

Las 10 categorías fueron seleccionadas mediante revisión de la literatura, consulta con profesionales y verificación contra currículos de finanzas empresariales.

\begin{table}[H]
\centering
\caption{Taxonomía de tareas financieras de FinanceBench}
\label{tab:taxonomia}
\begin{tabularx}{\textwidth}{clXXc}
\toprule
\textbf{\#} & \textbf{Categoría} & \textbf{Tipo de razonamiento} & \textbf{Dificultad típica} & \textbf{Casos} \\
\midrule
1 & Cálculo de VPN (sobre FCFE) & Aritmético + financiero & Media--Alta & 100 \\
2 & Cálculo de TIR (sobre FCFE) & Aritmético iterativo & Alta & 100 \\
3 & Proyección de Flujo de Caja & Aritmético secuencial & Media & 100 \\
4 & Generación de Escenarios & Mixto (aritmético + contextual) & Media--Alta & 100 \\
5 & Análisis de Punto de Equilibrio & Aritmético + conceptual & Media & 100 \\
6 & Evaluación de Estructura de Capital & Conceptual + aritmético & Media--Alta & 100 \\
7 & Cálculo de Indicadores Financieros & Aritmético + identificación & Media & 100 \\
8 & Interpretación de Estados Financieros & Interpretativo + numérico & Media & 100 \\
9 & Análisis de Sensibilidad & Aritmético repetitivo & Alta & 100 \\
10 & Recomendaciones de Viabilidad & Integrador & Alta & 100 \\
\bottomrule
\end{tabularx}
\smallskip
\footnotesize\textit{Fuente: Elaboración propia.}
\end{table}

Cada una de las 10 categorías evalúa un tipo de tarea financiera distinto, generando múltiples campos medibles por caso (ingresos, costos, depreciación, VPN, TIR, indicadores, etc.). En la Sección~\ref{sec:resultados}, los resultados se reportan tanto por categoría de tarea (rendimiento global en cada tipo de caso) como por tipo de métrica financiera (exactitud en campos específicos transversales a todos los casos), proporcionando una visión complementaria del rendimiento.

\subsection{Criterios de Dificultad}

La Tabla~\ref{tab:dificultad} define los criterios operacionales que determinan la clasificación de cada caso.

\begin{table}[H]
\centering
\caption{Criterios operacionales de dificultad}
\label{tab:dificultad}
\begin{tabularx}{\textwidth}{lXXX}
\toprule
\textbf{Criterio} & \textbf{Simple} & \textbf{Medio} & \textbf{Complejo} \\
\midrule
Fuentes de ingreso & 1 & 2--3 & 4+ \\
Estructura de costos & Fijos constantes, variables proporcionales & Costos escalonados & Costos interdependientes, no lineales \\
Financiamiento & Solo capital propio & 1 fuente de deuda & Mixto (2+ fuentes, tasas diferenciadas) \\
Depreciación & 1 activo, línea recta & 2 activos & 3+ activos con vidas útiles diferentes \\
Pérdidas fiscales & No aplica & Posibles, 1 año & Consecutivas, con arrastre multi-año \\
Horizonte & 5 años estándar & 5 años con valor residual & 5--7 años con reinversiones \\
Impuestos & Tasa fija 35\% & Tasa fija 35\% + deducciones & Regímenes diferenciados \\
\bottomrule
\end{tabularx}
\smallskip
\footnotesize\textit{Fuente: Elaboración propia.}
\end{table}

\subsection{Proceso de Construcción del Dataset}

\textbf{Paso 1 --- Generación de casos.} Cada categoría contiene 100 casos que cubren 10 sectores de negocio: cafetería, software/tecnología, retail/comercio, construcción, manufactura, consultoría, e-commerce, agricultura, bienes raíces y alojamiento. Los ejes de diversidad incluyen: inversión inicial (5--500 millones COP), estructura de costos, horizonte de proyección (3--7 años), tasa de descuento del accionista (8--18\%), y tasas de crecimiento de ingresos. Cada caso fue evaluado bajo cuatro escenarios económicos ---optimista, normal, pesimista y controlado---, generando $\sim$296 campos comparables por caso.

\textbf{Paso 2 --- Estructura estandarizada y cálculo de impuestos.} Todos los casos siguen un formato uniforme: (a) enunciado completo, (b) pregunta específica, (c) solución determinística, (d) nivel de dificultad, (e) categoría y sector.

La tasa de impuesto sobre la renta corporativa por defecto es del \textbf{35\%}, correspondiente a la tasa estatutaria de la Ley 2155 de 2021 (Ley de Inversión Social), que modificó el Artículo 240 del Estatuto Tributario de Colombia. El mecanismo de \textbf{pérdidas fiscales acumulables} sigue la lógica tributaria estándar: si la utilidad antes de impuestos es negativa en el año $t$, no se paga impuesto y la pérdida se arrastra para compensar utilidades futuras; esto refleja el tratamiento fiscal real del artículo 147 del Estatuto Tributario colombiano.


\begin{itemize}
  \item Si en el año $t$ la EBT es \textbf{negativa}: no se pagan impuestos; la pérdida acumulada se incrementa en $|\text{EBT}_t|$.
  \item Si en el año $t$ la EBT es \textbf{positiva}: $\text{base\_gravable} = \max(0,\; \text{EBT}_t - \text{pérdida\_acumulada})$; la pérdida acumulada se reduce en el monto aplicado.
  \item Las pérdidas \textbf{nunca se compensan contra nuevas pérdidas}, solo contra utilidades futuras.
\end{itemize}

Ejemplo del formato y derivación de FCFE:

\begin{quote}
\textbf{Caso ejemplo --- Cálculo de VPN (Cafetería, Simple)}

\textit{Enunciado:} Inversión total \$50.000.000 COP: \$30.000.000 capital propio + \$20.000.000 crédito bancario a 5 años, tasa 10\% anual, cuota fija. Depreciación anual \$3.000.000. Tasa de descuento del accionista: 15\%. Tasa de impuesto: 35\%.

\textit{Derivación del FCFE Año 1:}\\
Utilidad operativa (EBIT): \$22.000.000\\
$-$ Intereses año 1: \$2.000.000\\
$=$ EBT: \$20.000.000\\
$-$ Impuestos (35\%): \$7.000.000\\
$=$ Utilidad Neta: \$13.000.000\\
$+$ Depreciación: \$3.000.000\\
$-$ Pago de capital (principal): \$4.000.000\\
$=$ \textbf{FCFE Año 1: \$12.000.000}

\textit{VPN:} $= -30.000.000 + \displaystyle\sum_{t=1}^{5} \frac{\text{FCFE}_t}{(1{,}15)^t}$\\
La inversión descontada es solo el aporte de equity (\$30.000.000).
\end{quote}

Los valores oráculo fueron recalculados determinísticamente usando aritmética de punto flotante de doble precisión. Este proceso, junto con la auditoría técnica del benchmark, identificó y corrigió errores en la implementación original del mecanismo de pérdidas fiscales acumulables y actualizó la tasa impositiva al 35\% vigente. Estas correcciones se realizaron antes de la evaluación de los modelos.

\textbf{Paso 3 --- Validación por panel de expertos.} Se conformó un panel de 5 profesionales financieros: 2 contadores públicos certificados con experiencia en evaluación de proyectos, 2 analistas financieros con práctica en consultoría para PYMEs, y 1 docente universitario de finanzas con experiencia en investigación cuantitativa. La experiencia profesional media del panel fue de 8 años.

Cada experto recibió un subconjunto aleatorio de 200 casos y fue instruido para: (a) resolver el caso de forma independiente; (b) comparar su solución con la de referencia; (c) clasificar discrepancias; (d) evaluar claridad del enunciado. Cada evaluador dedicó aproximadamente 40 horas al proceso. Las discrepancias se resolvieron en 3 sesiones de consenso.

El coeficiente kappa de Cohen \citep{Cohen1960, Artstein2008} obtenido fue de $\kappa = 0{,}87$, correspondiente a un nivel de acuerdo \textit{casi perfecto} según la escala de Landis y Koch \citeyear{LandisKoch1977}. El proceso resultó en la modificación de 27 soluciones de referencia (2,7\% del total), principalmente en casos complejos con pérdidas fiscales consecutivas y estructuras de financiamiento mixtas.

\subsection{Modelos Evaluados}

\begin{table}[ht]
\centering
\caption{Modelos evaluados en FinanceBench}
\label{tab:modelos}
\begin{tabularx}{\linewidth}{lllX}
\toprule
\textbf{Modelo} & \textbf{Versión} & \textbf{Tipo} & \textbf{Proveedor} \\
\midrule
Claude Opus 4.6  & claude-opus-4-6   & Propietario    & Anthropic \\
Claude Sonnet 4.6 & claude-sonnet-4-6 & Propietario   & Anthropic \\
Claude Haiku 4   & claude-haiku-4    & Propietario    & Anthropic \\
GPT-5.4          & gpt-5.4           & Propietario    & OpenAI    \\
GPT-5.3          & gpt-5.3           & Propietario    & OpenAI    \\
Gemini 3.1 Pro   & gemini-3.1-pro    & Propietario    & Google    \\
Gemini 3 Flash   & gemini-3-flash    & Propietario    & Google    \\
DeepSeek Chat    & deepseek-chat     & Código abierto & DeepSeek  \\
Grok Code        & grok-code         & Código abierto & xAI       \\
\bottomrule
\multicolumn{4}{l}{\footnotesize\textit{Fuente: Elaboración propia.}}
\end{tabularx}
\end{table}

\subsection{Protocolo de Evaluación}

\textbf{Estrategia de prompting.} Prompt estandarizado zero-shot:

\begin{quote}
\textit{``Eres un analista financiero experto. Se te presenta el siguiente caso de análisis financiero. Lee cuidadosamente todos los datos proporcionados y responde con precisión numérica. Muestra todos los pasos de tu cálculo. [ENUNCIADO] [PREGUNTA]''}
\end{quote}

Es importante señalar que el prompt no instruye explícitamente al modelo sobre la convención FCFE. Esta decisión fue deliberada: buscamos medir si los modelos aplican la convención correcta dada la formulación del caso (donde los flujos se presentan como ``flujos de caja netos para el accionista'' y la inversión se presenta como el aporte de equity). Sin embargo, esto implica que un modelo que interprete el caso bajo la convención FCFF podría ser penalizado, una limitación reconocida en la Sección~\ref{sec:limitaciones}.

\textbf{Parámetros técnicos.} Temperatura: 0. Tokens máximos: 1.000. Cada caso evaluado 1 vez en baseline (9.000 evaluaciones totales); subconjunto de 200 casos con 3 ejecuciones a temperatura 0,3 para estimar varianza.

\textbf{Extracción de respuestas.} Pipeline automatizado con tolerancia de $\pm$1\%. Campos sin valor numérico clasificados como error ``nulo/ausente''.

\subsection{Métricas de Evaluación}

\textbf{Exactitud (\textit{Accuracy}).} Porcentaje de campos correctos ($\pm$1\%). Un modelo con 88,47\% acierta en $\sim$88 de cada 100 campos.

\textbf{MAPE (\textit{Mean Absolute Percentage Error}).} Magnitud promedio de las desviaciones en campos incorrectos. Ejemplo concreto: si el VPN correcto de un proyecto es \$50.000.000 COP y el modelo reporta \$46.085.000, la desviación es de 7,83\% --- el modelo se equivocó por \$3.915.000. Un MAPE de 41,57\% significa que, en promedio, los campos erróneos se desvían del valor real por más de dos quintas partes.

\textbf{Precisión Conceptual.} ¿Aplicó la fórmula o método correcto? Binario (sí/no).

\textbf{Tasa de Alucinación.} Porcentaje de casos con valores no derivables de los datos del enunciado.

\textbf{Completitud.} Porcentaje de campos solicitados que el modelo completó.

\textbf{Consistencia.} Coeficiente de variación en ejecuciones múltiples del mismo caso.

% ============================================================
\section{Resultados}\label{sec:resultados}
% ============================================================

Esta sección presenta los hallazgos cuantitativos del benchmark de manera ordenada, avanzando desde el panorama general hacia los análisis más granulares. Cada tabla y figura se acompaña de un párrafo interpretativo que traduce los números a implicaciones concretas. Todas las cifras provienen directamente de las 2,66 millones de comparaciones campo a campo entre las respuestas de los 9 modelos y los valores oráculo calculados determinísticamente bajo la convención FCFE descrita en la Sección~\ref{sec:metodologia}. La discusión profunda de las causas e implicaciones se reserva para la Sección~\ref{sec:discusion}; aquí nos concentramos en \textit{qué} encontramos.

\subsection{Rendimiento General por Modelo}

La Tabla~\ref{tab:ranking} presenta el ranking completo de los 9 modelos evaluados según exactitud global y MAPE. La Figura~\ref{fig:ranking} visualiza esta jerarquía de manera inmediata.

\begin{table}[ht]
\centering
\caption{Rendimiento general de los modelos evaluados en FinanceBench}
\label{tab:ranking}
\begin{tabularx}{\linewidth}{clXXl}
\toprule
\textbf{Rango} & \textbf{Modelo} & \textbf{Exactitud (\%)} & \textbf{MAPE (\%)} & \textbf{Nivel} \\
\midrule
1 & Claude Opus 4.6   & 88,47 & 7,83  & Superior        \\
2 & Claude Sonnet 4.6 & 86,31 & 10,27 & Superior        \\
3 & GPT-5.4           & 83,92 & 12,41 & Superior        \\
4 & Gemini 3.1 Pro    & 81,84 & 14,95 & Superior        \\
5 & GPT-5.3           & 78,26 & 18,73 & Intermedio      \\
6 & Claude Haiku 4    & 75,13 & 22,14 & Intermedio      \\
7 & DeepSeek Chat     & 69,71 & 27,86 & Intermedio-Bajo \\
8 & Grok Code         & 62,38 & 36,29 & Bajo            \\
9 & Gemini 3 Flash    & 57,64 & 41,57 & Bajo            \\
\bottomrule
\multicolumn{5}{l}{\footnotesize\textit{Fuente: Elaboración propia. Exactitud calculada sobre $\sim$296 campos $\times$ 1.000 casos $=$ $\sim$296.000 comparaciones por modelo.}}
\end{tabularx}
\end{table}

\begin{figure}[ht]
\centering
\rule{0.9\linewidth}{5.5cm}
\caption{Exactitud global (\%) de los 9 LLM evaluados en FinanceBench, ordenados de mayor a menor rendimiento. Se observa una estructura de tres niveles claramente definida: nivel superior ($>$81\%, cuatro modelos), nivel intermedio (69--79\%, tres modelos) y nivel bajo ($<$63\%, dos modelos). La brecha entre el primer y el último modelo asciende a 30,83 puntos porcentuales.}
\label{fig:ranking}
\end{figure}

Los resultados revelan una estructura de tres niveles que emerge de las agrupaciones naturales en la distribución de puntuaciones. El \textbf{nivel superior} (exactitud $>$ 81\%) agrupa a Claude Opus 4.6, Claude Sonnet 4.6, GPT-5.4 y Gemini 3.1 Pro, todos con MAPE inferior al 15\%. Dentro de este nivel, la dispersión es modesta: 6,63 puntos porcentuales separan al primero del cuarto. El \textbf{nivel intermedio} (exactitud 69--79\%) comprende a GPT-5.3, Claude Haiku 4 y DeepSeek Chat, con MAPE entre 18,73\% y 27,86\%. El \textbf{nivel bajo} (exactitud $<$ 63\%) contiene a Grok Code y Gemini 3 Flash, con MAPE superiores al 36\%.

La brecha de 30,83 puntos porcentuales entre el mejor y el peor modelo es notablemente mayor que las brechas observadas en benchmarks de generación de texto general (10--15 puntos), confirmando que el cálculo financiero es una capacidad de alta varianza entre LLM actuales. La correlación entre exactitud y MAPE es fuertemente inversa (Pearson $r = -0{,}98$, $p < 0{,}001$): cada punto porcentual de exactitud ganado corresponde a aproximadamente 1,10 puntos de MAPE reducido.

\subsection{Rendimiento por Categoría de Métrica Financiera}

No todas las tareas financieras son igualmente difíciles para los LLM. La Tabla~\ref{tab:metricas} presenta la exactitud desglosada por categoría de métrica para los cuatro modelos de nivel superior, y las Figuras~\ref{fig:metricas} y~\ref{fig:heatmap} visualizan este análisis para todos los modelos.

\begin{table}[ht]
\centering
\caption{Exactitud por categoría de métrica (modelos de nivel superior)}
\label{tab:metricas}
\begin{tabularx}{\linewidth}{lXXXX}
\toprule
\textbf{Categoría de métrica} & \textbf{Claude Opus 4.6} & \textbf{Claude Sonnet 4.6} & \textbf{GPT-5.4} & \textbf{Gemini 3.1 Pro} \\
\midrule
Depreciación             & 99,2\% & 98,8\% & 98,5\% & 97,9\% \\
Gastos fijos             & 98,7\% & 98,1\% & 97,4\% & 96,8\% \\
Gastos variables         & 94,3\% & 92,7\% & 90,1\% & 88,6\% \\
Proyecciones de ingresos & 91,8\% & 89,5\% & 87,2\% & 85,4\% \\
Punto de equilibrio      & 88,1\% & 85,9\% & 82,7\% & 80,3\% \\
Período de recuperación  & 84,6\% & 82,1\% & 79,8\% & 77,2\% \\
TIR (sobre FCFE)         & 78,3\% & 75,6\% & 72,4\% & 69,8\% \\
VPN (sobre FCFE)         & 76,9\% & 74,2\% & 70,8\% & 68,1\% \\
\bottomrule
\multicolumn{5}{l}{\footnotesize\textit{Fuente: Elaboración propia. VPN y TIR calculados sobre flujos de caja apalancados (FCFE).}}
\end{tabularx}
\end{table}

\begin{figure}[ht]
\centering
\rule{0.9\linewidth}{5cm}
\caption{Exactitud de los 9 LLM en cada una de las métricas financieras evaluadas. La depreciación y los gastos fijos rozan el rendimiento perfecto incluso en los modelos de menor capacidad, mientras que VPN y TIR presentan las mayores dificultades. El gradiente de dificultad es consistente entre modelos.}
\label{fig:metricas}
\end{figure}

\begin{figure}[ht]
\centering
\rule{0.9\linewidth}{5cm}
\caption{Mapa de calor que cruza las 10 métricas financieras (eje horizontal) con los 9 modelos evaluados (eje vertical). La codificación cromática va de rojo (baja exactitud) a verde (alta exactitud). Claude Opus 4.6 lidera en prácticamente todas las métricas. Los datos clave son: VPN oscila entre 48,6\% (Gemini 3 Flash) y 75,1\% (Claude Opus 4.6); Payback oscila entre 63,7\% y 97,3\%. La zona roja se concentra en VPN y TIR para los modelos de nivel bajo.}
\label{fig:heatmap}
\end{figure}

El patrón que emerge es revelador. Los cálculos de depreciación y gastos fijos rozan el rendimiento perfecto incluso en los modelos de menor capacidad (Gemini 3 Flash supera el 91\% en depreciación), porque involucran aritmética simple y repetitiva. En el extremo opuesto, el VPN y la TIR ---calculados sobre flujos FCFE que requieren cadenas de descuento compuesto de al menos seis operaciones dependientes--- caen al rango de 69--78\% para los mejores modelos, cuantificando con precisión el costo de la profundidad computacional. La correlación de rangos de Spearman entre cualquier par de modelos supera 0,94, lo que confirma que la dificultad es intrínseca al tipo de cálculo, no específica de un modelo.

\subsection{Rendimiento por Tipo de Proyecto}

El benchmark incluye 10 sectores de negocio: cafetería, software/tecnología, retail/comercio, construcción, manufactura, consultoría, e-commerce, agricultura, bienes raíces y alojamiento. La Figura~\ref{fig:sectores} presenta un mapa de calor que cruza modelos con sectores, permitiendo identificar si algún dominio amplifica o comprime las diferencias entre modelos.

\begin{figure}[ht]
\centering
\rule{0.9\linewidth}{5cm}
\caption{Mapa de calor de exactitud (\%) cruzando los 9 LLM (eje vertical) con los 10 sectores de negocio evaluados (eje horizontal). Los sectores con estructuras financieras más complejas (manufactura, bienes raíces) amplifican la diferencia entre modelos fuertes y débiles, mientras que los sectores más simples (cafetería, consultoría) la comprimen.}
\label{fig:sectores}
\end{figure}

Claude Opus 4.6 oscila entre 86,1\% (manufactura) y 91,2\% (cafetería), un rango de 5,1 puntos que refleja capacidad generalizada e independiente del sector. Gemini 3 Flash oscila entre 49,8\% y 64,7\%, un rango de 14,9 puntos que revela sensibilidad marcada a la complejidad del dominio. Los sectores con mayor complejidad financiera ---manufactura (múltiples líneas de producción, inventarios), bienes raíces (proyectos multi-fase, financiamiento complejo) y agricultura (períodos de maduración, estacionalidad)--- amplifican la diferencia entre modelos, mientras que sectores con estructuras más lineales (cafetería, consultoría) la comprimen parcialmente.

\subsection{Modelos Propietarios vs.\ Código Abierto}

La comparación entre modelos propietarios y de código abierto tiene alta relevancia práctica por sus implicaciones de accesibilidad económica. La Figura~\ref{fig:propietario} visualiza la distribución de rendimiento de ambos grupos.

\begin{figure}[ht]
\centering
\rule{0.9\linewidth}{5cm}
\caption{Comparación de la distribución de exactitud global entre los 7 modelos propietarios y los 2 modelos de código abierto. Los propietarios alcanzan una media de 78,80\% con amplia dispersión interna (57,64\%--88,47\%), mientras que los de código abierto promedian 66,05\% (62,38\%--69,71\%). La brecha de 12,75 puntos porcentuales es significativa en la práctica pero se concentra en las tareas de mayor complejidad computacional.}
\label{fig:propietario}
\end{figure}

Los siete modelos propietarios alcanzaron una exactitud media de 78,80\%, frente al 66,05\% de los dos modelos de código abierto, una brecha de 12,75 puntos porcentuales. Sin embargo, el grupo propietario es internamente heterogéneo: DeepSeek Chat (69,71\%) supera a Gemini 3 Flash (57,64\%) por 12,07 puntos. La brecha se comporta asimétricamente según la complejidad: en tareas simples se reduce a $\sim$7 puntos; en tareas complejas de flujo de caja descontado se amplía a 21,9 puntos. La prueba $t$ de Welch arroja $t(3{,}41) = 2{,}87$, $p = 0{,}054$, con significancia estadística marginal y tamaño del efecto grande ($d$ de Cohen $= 1{,}52$).

\subsection{Análisis de Errores}

La taxonomía de errores revela patrones cualitativamente distintos entre niveles de rendimiento. La Figura~\ref{fig:errores} presenta la distribución de tipos de error para cada modelo.

\begin{figure}[ht]
\centering
\rule{0.9\linewidth}{5cm}
\caption{Distribución porcentual de los 5 tipos de error (aritmético, conceptual, alucinación, formato, nulo/ausente) para cada modelo, en barras apiladas al 100\%. Los modelos de mayor capacidad concentran sus errores en la categoría aritmética, mientras que los de menor capacidad presentan mayor proporción de errores conceptuales y alucinaciones.}
\label{fig:errores}
\end{figure}

\begin{table}[ht]
\centering
\caption{Distribución de tipos de error (\% del total de errores por modelo)}
\label{tab:errores}
\begin{tabularx}{\linewidth}{lXXXX}
\toprule
\textbf{Tipo de error} & \textbf{Claude Opus 4.6} & \textbf{GPT-5.4} & \textbf{DeepSeek Chat} & \textbf{Gemini 3 Flash} \\
\midrule
Aritmético   & 47,2\% & 41,8\% & 33,6\% & 26,1\% \\
Conceptual   & 23,1\% & 28,4\% & 30,2\% & 31,8\% \\
Alucinación  & 11,4\% & 13,2\% & 16,8\% & 19,4\% \\
Formato      & 12,8\% & 10,1\% &  9,7\% &  8,9\% \\
Nulo/Ausente &  5,5\% &  6,5\% &  9,7\% & 13,8\% \\
\bottomrule
\multicolumn{5}{l}{\footnotesize\textit{Fuente: Elaboración propia.}}
\end{tabularx}
\end{table}

El hallazgo más revelador es el cambio composicional: los modelos de nivel superior cometen predominantemente errores aritméticos (47,2\% en Opus), indicando que cuando se equivocan, generalmente han identificado el método correcto pero fallan en la ejecución numérica. Los modelos de nivel bajo muestran proporción mayor de errores conceptuales (31,8\%) y alucinaciones (19,4\%), indicando carencias más fundamentales en comprensión financiera. Esta diferencia de composición explica por qué la brecha en MAPE excede proporcionalmente la brecha en exactitud: los errores conceptuales producen desviaciones de magnitud mucho mayor que los aritméticos.

La mediana del error porcentual absoluto condicionado confirma el patrón: 3,21\% para Claude Opus 4.6 versus 18,74\% para Gemini 3 Flash. El percentil 90 alcanza 12,8\% versus 67,3\% respectivamente, evidenciando fallos catastróficos ocasionales en los modelos inferiores.

DeepSeek Chat presenta un patrón notable: su tasa elevada de campos nulos (9,7\%), concentrada en VPN y TIR, sugiere una ``abstención calibrada'' ---el modelo declina responder en lugar de inventar---. Su exactitud sobre campos completados sube a $\sim$74,8\%, una estrategia conservadora preferible en contextos financieros.

\subsection{Efecto del Nivel de Dificultad}

La estratificación en tres niveles de dificultad permite evaluar la degradación ante exigencias crecientes. La Figura~\ref{fig:dificultad} visualiza estas trayectorias.

\begin{figure}[ht]
\centering
\rule{0.9\linewidth}{5cm}
\caption{Exactitud de los 9 LLM en los tres niveles de dificultad (Simple, Medio, Complejo). Todos los modelos muestran degradación progresiva, pero la tasa de degradación es inversamente proporcional a la capacidad del modelo: los de nivel superior pierden $\sim$5 puntos porcentuales; los de nivel bajo pierden $\sim$6,4 puntos. Esta degradación diferencial es un indicador de robustez.}
\label{fig:dificultad}
\end{figure}

\begin{table}[ht]
\centering
\caption{Exactitud por nivel de dificultad}
\label{tab:dificultad}
\begin{tabularx}{\linewidth}{lXXXX}
\toprule
\textbf{Modelo} & \textbf{Simple (\%)} & \textbf{Medio (\%)} & \textbf{Complejo (\%)} & \textbf{Brecha (pp)} \\
\midrule
Claude Opus 4.6   & 90,52 & 88,31 & 85,89 & 4,63 \\
Claude Sonnet 4.6 & 88,41 & 86,17 & 83,58 & 4,83 \\
GPT-5.4           & 86,14 & 83,78 & 81,07 & 5,07 \\
Gemini 3.1 Pro    & 84,03 & 81,62 & 79,14 & 4,89 \\
GPT-5.3           & 80,58 & 78,09 & 75,22 & 5,36 \\
Claude Haiku 4    & 77,49 & 75,01 & 71,84 & 5,65 \\
DeepSeek Chat     & 72,18 & 69,53 & 66,27 & 5,91 \\
Grok Code         & 65,02 & 62,14 & 58,87 & 6,15 \\
Gemini 3 Flash    & 60,34 & 57,38 & 53,92 & 6,42 \\
\bottomrule
\multicolumn{5}{l}{\footnotesize\textit{Fuente: Elaboración propia. Brecha = simple $-$ complejo en puntos porcentuales.}}
\end{tabularx}
\end{table}

Los modelos de nivel superior pierden entre 4,63 y 5,07 puntos al pasar de casos simples a complejos; los de nivel bajo pierden entre 6,15 y 6,42 puntos. Los casos complejos introducen múltiples fuentes de ingreso, depreciación de varios activos, estructuras de financiamiento mixtas y el mecanismo de pérdidas fiscales acumulables, amplificando los errores en cascada. Claude Opus 4.6 exhibe el menor factor de amplificación (1,8$\times$ en cinco años), Gemini 3 Flash el mayor (3,7$\times$).

\subsection{Rendimiento por Escenario Económico}

Cada caso fue evaluado bajo cuatro escenarios económicos. La Figura~\ref{fig:escenarios} presenta la exactitud desglosada por escenario para todos los modelos.

\begin{figure}[ht]
\centering
\rule{0.9\linewidth}{5cm}
\caption{Exactitud de los 9 LLM según el escenario económico evaluado (Optimista, Normal, Pesimista, Controlado). El escenario normal produce la mayor exactitud promedio (79,4\%), mientras que el pesimista produce la menor (72,1\%). Los modelos de nivel superior mantienen estabilidad entre escenarios (rango de $\sim$5 pp); los de nivel bajo exhiben variabilidad marcada (rango de hasta 18,3 pp).}
\label{fig:escenarios}
\end{figure}

El escenario normal produjo la mayor exactitud promedio (79,4\%), seguido del optimista (76,8\%), controlado (75,2\%) y pesimista (72,1\%). Claude Opus 4.6 mantuvo la mayor estabilidad entre escenarios (rango de 4,8 pp), mientras que Grok Code exhibió el rango más amplio (18,3 pp). El escenario pesimista es particularmente desafiante porque exige razonar simultáneamente sobre tasas de descuento más altas, ingresos reducidos y costos incrementados. La brecha entre mejor y peor modelo se amplía a 33,7 pp en el escenario pesimista (vs.\ 24,1 pp en el controlado), confirmando que las condiciones adversas funcionan como prueba de estrés que separa más agresivamente a los modelos por capacidad.

Esta observación tiene una implicación preocupante: los escenarios pesimistas son los más importantes para la toma de decisiones de inversión, y es precisamente allí donde la exactitud de los LLM es menor.

\subsection{Predicción de Viabilidad}

Más allá de la exactitud en métricas individuales, la aplicación práctica más relevante del benchmark para emprendedores es la clasificación binaria de viabilidad: dado el VPN calculado (sobre FCFE), ¿el modelo clasifica correctamente el proyecto como viable (VPN $>$ 0) o no viable (VPN $\leq$ 0)? La Figura~\ref{fig:viabilidad} presenta las matrices de confusión de esta tarea para cada modelo.

\begin{figure}[ht]
\centering
\rule{0.9\linewidth}{6cm}
\caption{Matrices de confusión para la tarea de clasificación binaria de viabilidad (Viable/No Viable basada en VPN $>$ 0) de los 9 LLM evaluados. F1-scores: Claude Opus 4.6 (0,914), Claude Sonnet 4.6 (0,886), GPT-5.4 (0,856), Gemini 3.1 Pro (0,830), GPT-5.3 (0,787), Claude Haiku 4 (0,749), DeepSeek Chat (0,687), Grok Code (0,634), Gemini 3 Flash (0,634). Los modelos de nivel superior clasifican correctamente más del 85\% de los proyectos.}
\label{fig:viabilidad}
\end{figure}

Claude Opus 4.6 alcanza un F1-score de 0,914 en la tarea de predicción de viabilidad, lo que significa que clasifica correctamente la decisión inversión/rechazo en más del 91\% de los casos. Este resultado es particularmente relevante porque, en la práctica, el emprendedor no necesita que el VPN sea exacto al centavo; necesita que la recomendación direccional (invertir o no invertir) sea correcta. Incluso modelos de nivel intermedio como Claude Haiku 4 (F1 = 0,749) y DeepSeek Chat (F1 = 0,687) alcanzan niveles de clasificación que, combinados con revisión humana en casos marginales, podrían ser útiles para filtrado inicial de proyectos.

El F1-score desciende de manera monotónica con el ranking general de exactitud, confirmando que la capacidad de cálculo numérico y la capacidad de clasificación de viabilidad están estrechamente acopladas en este dominio.

% ============================================================
\section{Análisis y Discusión}\label{sec:discusion}
% ============================================================

Si la Sección~\ref{sec:resultados} respondió \textit{qué encontramos}, esta sección aborda \textit{por qué} los resultados son como son, \textit{qué significan} y \textit{hacia dónde} apuntan.

\subsection{La Hipótesis de la Profundidad Computacional}

El hallazgo más consistente es el gradiente de dificultad entre tipos de métricas: la exactitud desciende predeciblemente a medida que el número de operaciones dependientes aumenta. La depreciación (una división y replicación) roza el 100\,\%. Las proyecciones de ingresos (multiplicaciones encadenadas) descienden al 85--92\,\%. El VPN y la TIR (cadenas de 10+ operaciones de descuento sobre FCFE) caen al 69--78\,\% incluso para los mejores modelos. La correlación de rangos de dificultad entre cualquier par de modelos supera 0,94, confirmando que estamos ante una propiedad del tipo de cálculo, no del modelo.

Proponemos la ``hipótesis de la profundidad computacional'': la exactitud de los LLM en cálculos financieros degrada de manera aproximadamente logarítmica con el número de pasos computacionales dependientes. Cada duplicación del número de operaciones encadenadas produce una disminución de 11--15 puntos porcentuales para modelos de nivel superior y 13--18 puntos para los de nivel inferior. Esta hipótesis es coherente con \citet{Dziri2023}, quienes demostraron en \textit{Faith and Fate} fallos sistemáticos de composicionalidad en LLM en tareas de razonamiento multi-paso.

La explicación técnica radica en la generación autoregresiva: el modelo no ``almacena'' un resultado intermedio como haría una hoja de cálculo, sino que lo genera como texto y lo ``relee'' cuando lo necesita. Cada transición generación-lectura introduce oportunidades de error que se propagan en cascada. Los datos cuantifican esta amplificación: factor de 1,8× en cinco años para Claude Opus 4.6 versus 3,7× para Gemini 3 Flash, lo que sugiere que Opus re-ancla sus cálculos en datos originales con mayor frecuencia.

La distribución bimodal del punto de equilibrio aporta evidencia complementaria. Esta tarea no exige alta profundidad computacional (una sola división) pero sí la selección correcta de una fórmula específica: $Q = \text{CF} / (P - \text{CV})$. Los modelos que la tienen internalizada aciertan consistentemente; los que no, cometen errores cualitativamente diferentes. Este comportamiento de umbral, contrastado con la degradación gradual del VPN, sugiere dos tipos de exigencia: selección del método (desafío de conocimiento, con umbral) y ejecución del cálculo (desafío aritmético, con degradación continua).

\subsection{Rendimiento General vs.\ Rendimiento Financiero}

Los modelos que lideran rankings generales ---Claude Opus 4.6, GPT-5.4--- también lideran en FinanceBench, lo que sugiere que una base sólida de capacidad general es condición necesaria para el rendimiento financiero. Sin embargo, el orden relativo no se preserva exactamente: Gemini 3.1 Pro compite cercanamente con GPT-5.4 en benchmarks generales pero queda 2,08 puntos por debajo en FinanceBench; Gemini 3 Flash obtiene resultados aceptables en tareas textuales pero colapsa al último lugar en cálculo financiero.

Estas discrepancias respaldan la tesis de que la capacidad general es necesaria pero no suficiente para el razonamiento financiero cuantitativo. El análisis financiero impone demandas que los benchmarks generales no capturan: precisión numérica sostenida, selección de fórmulas específicas, manejo simultáneo de escenarios e interpretación contextual. La distinción entre saber (conocimiento declarativo, lo que MMLU evalúa) y hacer (conocimiento procedimental, lo que FinanceBench evalúa) justifica la existencia de benchmarks especializados.

Respecto a la pregunta de investigación secundaria 4 (exploratoria): los datos sugieren que los rankings generales predicen parcialmente el rendimiento financiero ---los modelos de nivel superior en benchmarks generales tienden a ocupar posiciones de nivel superior en FinanceBench---, pero con reordenamientos significativos que impiden la predicción exacta. Un estudio de correlación formal entre scores de MMLU y FinanceBench, con una muestra ampliada de modelos, constituye una dirección prioritaria de trabajo futuro.

\subsection{Patrones de Fallo Sistemáticos}

Los \textbf{errores aritméticos en cascada} son el modo de fallo dominante para modelos de nivel superior. Un error del 2\,\% en la proyección de ingresos del año 1 se amplifica a un 5--8\,\% en el año 5, produciendo proyecciones materialmente incorrectas que, sin embargo, parecen razonables a simple vista.

La \textbf{contaminación entre escenarios} ocurre cuando un modelo hereda valores del escenario normal al generar el pesimista: 3,2\,\% de casos en Gemini 3 Flash, 2,1\,\% en Grok Code, menos de 0,5\,\% en nivel superior.

La \textbf{mala aplicación de la tasa de descuento} es el error individual de mayor impacto por magnitud: puede generar desviaciones en VPN superiores al 100\,\%. Representa el 34\,\% de todos los errores de VPN que exceden el 50\,\% de desviación. Claude Opus 4.6 demuestra 99,1\,\% de precisión en la tasa de descuento como campo aislado; Gemini 3 Flash desciende a 87,3\,\%.

Tres factores convergen: una \textbf{limitación arquitectónica} (la generación token a token no fue diseñada para precisión numérica sostenida), una \textbf{insuficiencia en datos de entrenamiento} (escasez de derivaciones financieras paso a paso, en contraste con abundancia de texto financiero descriptivo), y la \textbf{ausencia de herramientas integradas} (a diferencia de un analista que usa una hoja de cálculo para los cálculos y reserva su juicio para la interpretación).

\subsection{Implicaciones Prácticas para Desarrolladores}

Para \textbf{proyecciones de alto impacto} (créditos, presentaciones a inversores, regulación): solo modelos de nivel superior, con verificación determinística de VPN y TIR. Arquitectura recomendada: el LLM genera supuestos y contextualización; un módulo determinístico ejecuta las fórmulas exactas sobre FCFE.

Para \textbf{análisis exploratorio} (evaluaciones preliminares, comparaciones rápidas): modelos de nivel intermedio ofrecen relación favorable entre precisión y costo. Su MAPE de 18--28\,\% es aceptable para orientación direccional. Claude Haiku 4 (75,13\,\%, MAPE 22,14\,\%) emerge como opción interesante por su bajo costo operativo.

Para \textbf{contextos educativos}: incluso modelos de nivel bajo pueden cumplir función útil si sus respuestas se presentan como ilustrativas.

\subsection{Implicaciones para Emprendedores y Usuarios Finales}

La regla general: \textbf{cuanto más simple sea el cálculo, más confiable es el resultado}. Márgenes y depreciación: $>$97\,\% de probabilidad de acierto. VPN y TIR: 70--78\,\%. Casos complejos con financiamiento mixto: probabilidad de error adicional.

Los errores más peligrosos no son los aritméticos (desviaciones pequeñas, detectables con intuición del negocio) sino los conceptuales (omitir el escudo fiscal en el WACC, usar la convención de flujos equivocada) y las alucinaciones (tasas o subsidios inventados).

Tres pautas de uso responsable: (1) verificar siempre VPN, TIR y WACC con herramienta independiente; (2) desconfiar de la confianza del modelo ---un párrafo bien redactado puede contener un número falso---; (3) usar el LLM para lo que hace bien: generar supuestos, estructurar análisis, identificar variables clave.

\subsection{Comparación con Estudios Previos}

Nuestros hallazgos dialogan productivamente con la literatura. La jerarquía de tres niveles es consistente con \citet{Xie2024}, que también reporta liderazgo de modelos de la familia GPT-4 y Claude en tareas financieras, aunque FinBen evalúa tareas predominantemente textuales donde la exactitud supera el 90\,\%. Nuestro benchmark, al evaluar la generación posterior de estos modelos (GPT-5.4, Claude Opus 4.6) en tareas cuantitativas abiertas, confirma y extiende ese hallazgo con un techo de 88,47\,\% ---consistente con la hipótesis de que el cálculo abierto impone una demanda adicional que la clasificación no captura.

La brecha propietario--código abierto de 12,75 puntos supera la de 8--10 puntos típicamente reportada en MMLU y GSM8K, coherente con nuestra hipótesis de profundidad computacional: las capacidades donde la brecha se amplifica son precisamente aquellas en las que los modelos propietarios de frontera han invertido más.

Tres contribuciones novedosas respecto a la literatura: la taxonomía de errores condicionada por capacidad del modelo; la cuantificación precisa del gradiente de dificultad entre métricas; y la evaluación en contexto latinoamericano.

La adopción de la convención FCFE implica que comparaciones directas con benchmarks que usen WACC + FCFF requieren cautela. Un LLM que responda consistentemente con FCFF produciría valores diferentes a nuestros oráculo sin que ello constituya error financiero. Esta observación refuerza la importancia de que futuros benchmarks declaren explícitamente su convención de flujos, un estándar de transparencia que FinanceBench adopta.

\subsection{Direcciones Futuras de Investigación}

\textbf{Evaluación con prompting avanzado.} Chain-of-thought, few-shot con ejemplos financieros y prompting estructurado pueden mejorar sustancialmente la precisión \citep{Wei2022,Sahoo2024}. FinanceBench permite esta comparación directamente.

\textbf{Arquitecturas híbridas.} Un sistema donde el LLM genera supuestos y un motor determinístico ejecuta las fórmulas sobre FCFE podría superar el 95\,\% de exactitud, eliminando la mayor parte de los errores aritméticos.

\textbf{Extensión multilingüe.} Evaluar si el rendimiento difiere entre enunciados en español e inglés, controlando por equivalencia de contenido.

\textbf{Fine-tuning especializado.} Entrenar modelos sobre derivaciones financieras paso a paso, no solo texto financiero descriptivo.

\textbf{Correlación formal con benchmarks generales.} Estudio con muestra ampliada de modelos que calcule la correlación entre scores de MMLU y FinanceBench, respondiendo plenamente la pregunta de investigación 4.

\textbf{Actualización periódica.} FinanceBench fue diseñado desde el inicio para ser extensible: nuevos modelos, categorías y niveles de dificultad pueden incorporarse sin alterar la infraestructura.

% ============================================================
\section{Limitaciones}\label{sec:limitaciones}
% ============================================================

Reconocer las limitaciones fortalece el paper porque demuestra que los investigadores comprenden el alcance real de sus conclusiones.

La primera limitación concierne al \textbf{alcance del benchmark}. Las 10 categorías cubren el núcleo del análisis de viabilidad, pero quedan fuera la valoración de empresas a largo plazo, el análisis de riesgo crediticio, la optimización de portafolios y la modelación de derivados. Esta delimitación fue deliberada: nos concentramos en las tareas más relevantes para emprendedores y PYMEs.

La segunda limitación se refiere a la \textbf{convención de flujos de caja}. FinanceBench calcula VPN y TIR sobre flujos de caja apalancados para el accionista (FCFE), descontados a la tasa del accionista ($k_e$). Un LLM que responda con la convención alternativa ---FCFF descontado a WACC--- producirá valores numéricamente diferentes sin que ello constituya error financiero. Adicionalmente, el prompt estándar no instruye explícitamente al modelo sobre la convención FCFE; aunque los enunciados formulan los flujos como ``flujos de caja netos para el accionista'' y la inversión como ``aporte de equity'', un modelo podría interpretar el caso bajo FCFF y ser penalizado. Futuros trabajos podrían evaluar ambas convenciones en paralelo.

La tercera limitación se refiere a la \textbf{naturaleza sintética de los casos}. Todos los 1.000 casos fueron construidos para el benchmark. Los casos sintéticos pueden ser ``demasiado limpios'' comparados con la realidad financiera, donde los datos son incompletos y los supuestos implícitos.

La cuarta involucra al \textbf{panel de validación}. Aunque el panel de 5 expertos logró un kappa de Cohen de 0,87 (acuerdo casi perfecto), su tamaño es modesto y sus perfiles presentan concentración geográfica en la costa caribe colombiana. Durante la auditoría técnica se identificaron y corrigieron errores en el mecanismo de pérdidas fiscales acumulables (\textit{tax loss carryforward}) para proyectos con pérdidas consecutivas, así como la actualización de la tasa impositiva al 35\,\% vigente (Ley 2155 de 2021). Estas correcciones, realizadas antes de la evaluación de modelos, constituyen evidencia del rigor del aseguramiento de calidad.

La quinta es la \textbf{estrategia de prompting}. Evaluamos con prompt zero-shot estándar, midiendo capacidad base. El rendimiento de todos los modelos probablemente mejore con prompting optimizado \citep{Wei2022,Sahoo2024}, pero la magnitud de la mejora varía entre modelos y categorías. FinanceBench permite esta comparación directamente.

La sexta se relaciona con el \textbf{idioma}. La mayoría de los LLM han sido entrenados predominantemente con texto en inglés, lo que podría afectar el rendimiento cuando los enunciados se presentan en español.

La séptima es la \textbf{velocidad de evolución del campo}. Los modelos evaluados corresponden a versiones específicas; versiones posteriores podrían rendir diferente. La diferencia de 5,66 puntos entre GPT-5.3 y GPT-5.4 ilustra la magnitud del cambio entre versiones. El diseño extensible del benchmark permite actualizaciones periódicas.

Finalmente, la \textbf{ausencia de evaluación cualitativa profunda}. FinanceBench mide si los números son correctos, pero no evalúa la calidad del razonamiento narrativo, la pertinencia de las advertencias sobre riesgos, ni la utilidad práctica de las recomendaciones. Estas dimensiones requieren juicio humano no escalable a 2,66 millones de comparaciones.

Cada una de estas limitaciones acota las conclusiones sin invalidarlas. Los hallazgos centrales son robustos dentro del marco definido.

% ============================================================
\section{Conclusiones}\label{sec:conclusiones}
% ============================================================

Este estudio nació de una observación preocupante: millones de personas utilizan LLM para tomar decisiones financieras sin que nadie haya medido sistemáticamente qué tan confiables son estos sistemas. FinanceBench fue construido para cerrar esa brecha.

Los cinco objetivos de investigación declarados en la Sección~\ref{sec:objetivos} fueron cumplidos. El \textbf{Objetivo 1} se alcanzó con el diseño de una taxonomía de 10 categorías de tareas financieras (Sección~\ref{sec:metodologia}). El \textbf{Objetivo 2} se cumplió mediante la definición de 6 métricas complementarias (Sección~\ref{sec:metodologia}). El \textbf{Objetivo 3} se materializó en un dataset de 1.000 casos validados por un panel de 5 expertos con kappa de 0,87 (Sección~\ref{sec:metodologia}). El \textbf{Objetivo 4} se ejecutó evaluando 9 LLM bajo condiciones estandarizadas (Sección~\ref{sec:resultados}). El \textbf{Objetivo 5} se abordó con la clasificación de errores en 5 tipos y el análisis de patrones por modelo y categoría (Secciones~\ref{sec:resultados} y~\ref{sec:discusion}).

Los hallazgos fundamentales se sintetizan en cuatro conclusiones. Primera, existe una jerarquía de rendimiento robusta: Claude Opus 4.6 lidera con 88,47\,\% de exactitud y 7,83\,\% de MAPE, mientras que Gemini 3 Flash cierra con 57,64\,\% y MAPE de 41,57\,\%, una brecha de 30,83 puntos porcentuales que confirma que la selección del modelo tiene consecuencias cuantitativas enormes. Segunda, la dificultad sigue un gradiente predecible determinado por la profundidad computacional: depreciación roza la perfección ($>$97\,\%); VPN y TIR se mantienen como fronteras de capacidad ($<$79\,\%). Tercera, la composición de errores cambia cualitativamente: modelos de nivel superior cometen predominantemente errores aritméticos menores (47,2\,\%), mientras que los de nivel inferior acumulan errores conceptuales (31,8\,\%) y alucinaciones (19,4\,\%). Cuarta, los modelos propietarios mantienen una ventaja de 12,75 puntos sobre los de código abierto, concentrada en tareas de alta complejidad computacional.

En la tarea de clasificación binaria de viabilidad ---la aplicación más directa para emprendedores---, Claude Opus 4.6 alcanza un F1-score de 0,914, clasificando correctamente la decisión inversión/rechazo en más del 91\,\% de los casos. Este resultado, sumado a las cinco direcciones de investigación futura detalladas en la Sección~\ref{sec:discusion}, constituye la hoja de ruta natural para el avance del campo.

La contribución de FinanceBench es triple: como recurso de evaluación (primer benchmark integral de análisis financiero cuantitativo con 1.000 casos validados), como evidencia empírica (línea base contra la cual medir avances futuros), y como guía práctica (recomendaciones diferenciadas por nivel de riesgo). Invitamos a la comunidad a utilizar, citar y extender FinanceBench. El dataset, código y resultados están disponibles en [URL del repositorio].

Cerramos con una reflexión. Vivimos un momento singular: por primera vez, sistemas computacionales producen análisis financieros que, en apariencia, son indistinguibles de los de un profesional humano. La oportunidad reside en la democratización: un emprendedor en una zona rural de Colombia podría acceder a una evaluación de calidad comparable a la de una \textit{startup} en Silicon Valley. La responsabilidad reside en la honestidad: esa evaluación solo tiene valor si es correcta, y nuestros datos muestran que la probabilidad de error puede oscilar entre el 12\,\% y el 42\,\% dependiendo del modelo. La investigación rigurosa es la herramienta que permite distinguir la oportunidad genuina del optimismo injustificado. FinanceBench es nuestra contribución a esa tarea.

% ============================================================
\section*{Declaración de Ética}
% ============================================================

\textbf{Transparencia.} El dataset, código y resultados se publican como recurso de acceso abierto, sin restricciones para investigación. La reproducibilidad completa es un requisito ético.

\textbf{Advertencia.} Ningún LLM evaluado ---incluido el de mayor rendimiento--- está validado para toma autónoma de decisiones financieras sin supervisión humana. Un rendimiento del 88\,\% en un benchmark estandarizado no garantiza rendimiento equivalente en escenarios reales. La IA debe complementar el juicio humano, no reemplazarlo.

\textbf{Ausencia de conflictos de interés.} Los autores no mantienen relaciones comerciales con ningún proveedor de modelos evaluados. La evaluación se realizó de manera independiente.

\textbf{Uso responsable.} FinanceBench fue creado para informar, no para legitimar. Instamos a los desarrolladores a comunicar las limitaciones de los modelos a sus usuarios, y a los emprendedores a tratar los resultados de IA como puntos de partida, no como conclusiones definitivas.

\textbf{Privacidad.} Todos los casos son sintéticos. No se usaron datos financieros reales de personas ni empresas.

% ============================================================
\bibliography{references}
% ============================================================

\end{document}
